% Modernised reading — fonds Alexandre Grothendieck, shelfmark 161-3
% (pages 1-54, the whole folder).
%
% This is an INTERPRETATION, not a transcription. It re-reads the pages in
% today's notation and vocabulary, states the hypotheses the manuscript leaves
% implicit, and drops the critical apparatus so that the mathematics reads
% straight through. Everything the brackets and underlines carried — a doubtful
% reading, a notation he used differently, a missing passage — moves into a
% footnote.
%
% It is held to being mathematically correct as it stands, not merely faithful
% to the page. Where the manuscript is loose or mistaken, the true statement is
% what appears here, and the footnote says what the page has. In any dispute of
% substance the transcription governs: whoever cites Grothendieck must cite
% batch-0{1,2,3}.fr.tex.
%
% Pass: Opus 5 (claude-opus-5), 2026-08-16 — first pass, unchecked against the
% pages by a human. The folder was read whole in one pass, all three batches
% (pages 1-54) together, and this is a SINGLE document for the whole folder
% rather than one file per batch: 161-3 is a dossier rassemblé whose threads
% cross the batch boundaries in both directions — page 41 completes a sentence
% begun on page 39, page 48 answers a question left open on page 40 — and the
% pages are followed here by thread, not in the order the leaves are bound.
\documentclass[11pt,a4paper]{article}
% Le préambule commun à toutes les transcriptions du fonds Grothendieck.
%
% Il tient en une idée : **l'incertitude se compose, elle ne se résout pas.**
% Une transcription qui lisse un mot douteux a détruit la seule chose pour
% laquelle on la faisait — un lecteur doit pouvoir savoir, à chaque endroit,
% ce qui était sur le feuillet et ce qui a été ajouté. Les six macros qui
% suivent sont donc l'appareil critique, et non de la décoration.
%
% Les mêmes commandes sont reconnues par `scripts/render.mjs`, qui produit la
% vue de lecture du site. Une macro ajoutée ici doit l'être aussi là-bas,
% faute de quoi le rendu échouera bruyamment — ce qui est voulu : un rendu
% silencieusement faux est pire qu'un rendu absent.

\ProvidesPackage{grothendieck}[2026/08/08 Transcriptions du fonds Grothendieck]

\usepackage{amsmath,amssymb,amsthm}
\usepackage{mathrsfs}
\usepackage{stmaryrd}
% Les diagrammes commutatifs sont partout dans ce fonds : c'est la langue
% même de la géométrie algébrique de Grothendieck, et une transcription qui
% les rendrait en prose ne transcrirait rien.
\usepackage{tikz-cd}
% Français par défaut — c'est la langue des feuillets — mais l'anglais est
% chargé aussi : la traduction et le résumé sont des documents anglais, et la
% typographie française (espaces avant les deux-points) y serait une faute.
% Ces documents basculent par \selectlanguage{english} en tête de corps.
\usepackage[english,french]{babel}
\usepackage{fontspec}
\usepackage{geometry}
\geometry{a4paper,margin=25mm,marginparwidth=28mm}
\usepackage{marginnote}
\usepackage{xcolor}
% ulem, not soul. `soul`'s \st and \ul are built on a character-by-character
% scan that XeLaTeX's font machinery breaks: the document fails with an
% undefined control sequence rather than degrading. ulem does the same two jobs
% and is Unicode-engine safe. `normalem` stops it hijacking \emph, which the
% transcriptions use for Grothendieck's own emphasis.
\usepackage[normalem]{ulem}
\usepackage{hyperref}
\hypersetup{colorlinks=true,linkcolor=black,urlcolor=black,pdfborder={0 0 0}}

% --- Métadonnées du lot -----------------------------------------------------
% Reprises telles quelles par le script de rendu, qui les affiche en tête de
% la vue de lecture. Elles fixent ce que le fichier prétend couvrir : sans
% elles, une transcription flotte sans point d'attache dans un fonds de seize
% mille feuillets.

\newcommand{\folder}[1]{\gdef\grothFolder{#1}}
\newcommand{\batch}[1]{\gdef\grothBatch{#1}}
\newcommand{\leaves}[2]{\gdef\grothFirst{#1}\gdef\grothLast{#2}}
\newcommand{\foldertitle}[1]{\gdef\grothFoldertitle{#1}}
\gdef\grothFolder{?}\gdef\grothBatch{?}\gdef\grothFirst{?}\gdef\grothLast{?}\gdef\grothFoldertitle{}

% --- Le résumé --------------------------------------------------------------
%
% Il ouvre la lecture modernisée, sous le titre « Résumé », et il est composé
% en petit. Ce n'est pas un ornement : c'est le seul passage du document qui
% ne soit pas des mathématiques, et un lecteur doit pouvoir le reconnaître
% avant de l'avoir lu — puis le sauter s'il connaît déjà le sujet, ou s'y
% arrêter s'il ne le connaît pas. Le filet qui le suit marque la reprise du
% corps.

\newenvironment{resume}
  {\small\setlength{\parskip}{0.45em}}
  {\par\medskip\hrule\medskip}

% --- Mots-clés ---------------------------------------------------------------
%
% En anglais, en clôture du résumé de la lecture modernisée : le vocabulaire
% moderne sous lequel chercher ce que les feuillets construisent. Le manifeste
% du site (`npm run manifest`) les extrait pour étiqueter la cote entière —
% cette ligne est la source unique des tags, il n'y a pas de fichier à côté.
\newcommand{\keywords}[1]{\par\smallskip\noindent
  {\footnotesize\textsc{Keywords} --- \textit{#1}}\par}

% --- L'appareil critique ----------------------------------------------------

% Le numéro de feuillet, en marge. C'est l'ancre qui permet de régler un
% désaccord sur une formule en regardant *un* feuillet plutôt qu'une chemise
% de six cents. Le site s'en sert aussi pour tourner les pages du fac-similé
% quand on fait défiler la transcription.
\newcommand{\leaf}[1]{%
  \marginnote{\footnotesize\color{gray!70}\textsf{#1}}%
  \label{leaf:#1}\ignorespaces}

% Illisible : on ne devine pas. Un mot inventé qui se lit comme les autres est
% le pire résultat possible.
\newcommand{\ill}{\textcolor{red!60!black}{[\,\ldots\,]}}

% Lecture douteuse : le mot est donné, et signalé comme incertain.
\newcommand{\uncertain}[1]{\uline{#1}}

% Ajout de l'éditeur — une lettre manquante, un mot sous-entendu.
\newcommand{\add}[1]{\textcolor{blue!50!black}{[#1]}}

% Biffé par Grothendieck lui-même. Ce qu'il a rayé fait partie du document :
% une rature dit souvent où la pensée a changé de direction.
\newcommand{\struck}[1]{\sout{#1}}

% Note du transcripteur, sur l'état matériel du feuillet.
\newcommand{\note}[1]{\par\smallskip{\small\color{gray!60!black}\textit{[#1]}}\par\smallskip}

% Note marginale de Grothendieck, distincte du corps du texte.
\newcommand{\marginal}[1]{\par\smallskip\noindent\hspace{1em}%
  {\small\color{gray!60!black}#1}\par\smallskip}

% --- Titre ------------------------------------------------------------------

\AtBeginDocument{%
  \begin{center}
    {\large\bfseries Fonds Alexandre Grothendieck --- cote n\textsuperscript{o}~\grothFolder}\\[2pt]
    {\itshape\grothFoldertitle}\\[4pt]
    {\small feuillets \grothFirst--\grothLast\ (lot \grothBatch)}\\[2pt]
    {\footnotesize\color{gray!60!black}Transcription établie d'après le fac-similé numérisé
     par l'Université de Montpellier.\\
     Les lectures incertaines sont soulignées, les illisibles notées [\,\ldots\,],
     les ajouts entre crochets.}
  \end{center}
  \vspace{1em}}

\endinput


\folder{161-3}
\batch{1}
\pages{1}{54}
\dating{[vers 1963-1973]}
\watermark{Édition de démonstration\\interprétation personnelle de l'œuvre}
\foldertitle{Topos : notes manuscrites (s.d.) — lecture modernisée du dossier entier}

\begin{document}

\section*{Résumé}

\begin{resume}
Une chemise portant au crayon le seul mot « Topos », cinquante-quatre pages
sans date, et — c'est la première chose à dire — \emph{pas un texte}. Ce
dossier a été rassemblé, non écrit : cinq chantiers sans rapport entre eux y
voisinent parce qu'ils touchent tous au même mot. Qui y cherche un argument
continu en inventera un. Ce qu'on peut faire, en revanche, c'est démêler les
fils, et c'est ce que cette lecture fait : elle suit chaque chantier d'un bout
à l'autre, dans l'ordre de son propre développement et non dans celui des
feuillets.

Ce désordre n'est pas une figure de style. Deux fils traversent les coupures :
la page 41 achève une phrase commencée page 39, et la page 48 répond à une
question posée page 40. Et surtout, à partir de la page 34, \emph{deux
manuscrits distincts sont reliés en alternance} — les pages paires portent des
notes de cours, les impaires une théorie des espèces de structure, l'alternance
étant exacte jusqu'à la page 48. Deux feuillets sont même reliés à l'envers :
la page 54 précède la page 53, ce que la numérotation des théorèmes et une
phrase coupée en deux établissent l'une et l'autre. Lu dans l'ordre des pages,
ce dossier paraît une suite de fragments interrompus ; lu par fils, il n'a
presque pas de lacunes.


\textbf{Multiplier deux espaces.} Le premier chantier (pages 3 à 14) part
d'une question qu'on peut poser sans rien savoir. Pour deux espaces
ordinaires, le produit est le produit cartésien, et personne n'y regarde à deux
fois. Mais un topos ne se donne pas par ses points : il se donne par ce qu'on
peut y poser comme données locales — ses faisceaux — et le produit doit donc
être l'endroit où vivent les données à \emph{deux} variables, locales
séparément en chacune. Grothendieck montre que cet objet est bien un topos et
bien le produit ; puis vient la surprise, qui occupe quatre pages : le produit
de deux espaces ne porte pas toujours le produit des topos. C'est l'écart
familier en algèbre linéaire entre les fonctions de deux variables et celles
qu'on obtient à partir des produits $f(x)g(y)$ — savoir une chose séparément en
$x$ et en $y$ n'est pas la savoir en $(x,y)$. Il isole exactement la condition
qui referme l'écart, et sa démonstration tient en trois passages à la borne
supérieure. Deux notes techniques ferment le chantier (pages 15 à 21).

\textbf{Quand la topologie est une combinatoire.} Le deuxième chantier
(pages 23 à 28) est le plus joli du dossier. Autour d'un point, en général, il
y a une infinité de voisinages de plus en plus petits, sans plus petit. Mais il
arrive qu'il y en ait un — et alors toute la topologie tient dans une relation
d'ordre, et réciproquement tout ensemble ordonné donne un tel espace. Topologie
et combinatoire coïncident exactement. Grothendieck monte le dictionnaire dans
les deux sens et en tire la notion de \emph{point essentiel}, dont dépendra la
suite.

\textbf{Une page d'algèbre égarée.} Le troisième chantier (pages 30 à 32) n'a
rien à faire là : deux façons de décrire la même chose lorsqu'une algèbre
ressemble localement à une algèbre de matrices. C'est le début de ce qu'on
appelle aujourd'hui le groupe de Brauer d'un topos, et on y voit Grothendieck
écrire le mot « gerbe », puis le barrer.

\textbf{Peut-on axiomatiser les corps ?} Le quatrième chantier (pages 35 à 47,
pages impaires) prend les théories elles-mêmes pour objets. Un anneau se
définit par des opérations et des identités, et cette définition a la propriété
d'être aveugle au monde où on l'interprète. Un corps, lui, se définit par une
condition d'un autre genre — « tout élément non nul est inversible » — qui parle
d'éléments un par un, et rien ne dit d'avance qu'elle se transporte. La réponse
n'est ni oui ni non mais une ligne de partage : avec les seuls moyens qui
définissent les anneaux, non ; dès qu'on s'autorise en plus la localisation,
oui — et l'axiome qui apparaît alors est exactement celui qu'on écrit
aujourd'hui pour un corps dans un topos. Les dernières pages du fil poussent la
machinerie jusqu'à un énoncé qu'on reconnaît : une théorie et sa catégorie de
modèles se déterminent l'une l'autre.

\textbf{Le cours.} Le cinquième chantier (pages 34 à 54, pages paires puis
toutes) est le seul endroit du dossier où Grothendieck dise pourquoi il fait ce
qu'il fait. Pourquoi remplacer les espaces par les topos ? Parce que, écrit-il,
c'est le cadre naturel de presque toute structure algébrique ; parce qu'on peut
y construire, pour chaque espèce de structure, un topos qui la \emph{classifie}
— « algèbre $=$ topologie ! », avec le point d'exclamation ; parce que ces
objets se recollent, se quotientent et se localisent là où les espaces
refusent. La géométrie algébrique, qui est le motif historique, vient en
dernier dans sa liste. Suit une énumération des morphismes de topos par
l'exemple, et deux pages qu'il faut signaler à qui ne lira rien d'autre : pour
un groupe $G$ et un sous-groupe $H$, se restreindre à $H$ et se localiser
au-dessus de $G/H$ sont la même opération — et si $G$ est le groupe fondamental
d'un espace, cet énoncé formel \emph{est} la théorie de Galois des revêtements,
qui tombe comme cas particulier d'une remarque sur les topos. Le dossier
s'achève sur trois pages de sites — topologies sur une catégorie, faisceau
associé, théorèmes de Giraud — qui sont, à la compression près, ce que SGA 4
mettra plusieurs exposés à établir.


Les noms modernes sous lesquels chercher : 2-produit et 2-limite de topos,
produit de locales et compacité locale, morphisme et point essentiels,
topologie d'Alexandrov, espace sobre et sobrification, topos classifiant,
dualité de Gabriel--Ulmer et catégories localement présentables, objet corps et
topos de Zariski, théorème de Giraud et foncteur faisceau associé, algèbres
d'Azumaya et groupe de Brauer.

\keywords{product of topoi, locale product, Alexandrov topology, essential point, sober space, classifying topos, Gabriel–Ulmer duality, Giraud's theorem, Azumaya algebra}
\end{resume}


\pagerange{3}{3}
\section*{Le produit de deux topos (pages 3 à 8)}

Dans tout ce qui suit, $C$ et $D$ sont de petits sites, $\widehat{C}$ la catégorie
des préfaisceaux d'ensembles sur $C$ et $\widetilde{C}$ celle des faisceaux ;
$i_{C} : \widetilde{C} \to \widehat{C}$ est l'inclusion, $a_{C}$ son adjoint à
gauche, le foncteur « faisceau associé », qui est exact à
gauche.\footnote{Cette distinction entre $\widehat{C}$ et $\widetilde{C}$ est
celle du manuscrit, constante dans tout le dossier, et rien n'y est plus
important : les pages 53 et 54 portent précisément sur l'inclusion
$\widetilde{C} \subset \widehat{C}$ et son adjoint à gauche exact. Nous ne les
confondons nulle part.} Les topos sont les topos de Grothendieck ; $E$, $F$
en désignent deux.

\pagerange{3}{3}
\subsection*{Quatre descriptions du même objet}

La page 3 pose, sans le nommer autrement que par la lettre, l'objet
\[
\Pi(E,F) \;=\; \bigl\{\, \varphi : E^{\circ} \to F \ \text{commutant aux
limites} \,\bigr\},
\]
et en donne aussitôt quatre lectures équivalentes :

\begin{itemize}
\item les foncteurs $E^{\circ} \to F$ commutant aux limites — c'est-à-dire
transformant les colimites de $E$ en limites de $F$ — ou, ce qui revient au
même, admettant un adjoint à gauche ;\footnote{L'équivalence des deux
formulations n'est pas formelle : un foncteur qui commute aux limites n'a
d'adjoint à gauche que sous une hypothèse d'engendrement et de taille. Elle
est ici acquise parce qu'un topos est une catégorie localement présentable,
donc cocomplète et à générateur petit ; c'est le théorème de l'adjoint spécial.
La page ne le dit pas.}
\item symétriquement, les foncteurs $F^{\circ} \to E$ commutant aux limites ;
\item les foncteurs $E^{\circ} \times F^{\circ} \to (\mathrm{Ens})$ commutant
aux limites \emph{en chaque variable séparément} ;
\item et, si $E \simeq \widetilde{C}$ et $F \simeq \widetilde{D}$, les
foncteurs $C^{\circ} \times D^{\circ} \to (\mathrm{Ens})$ qui sont des faisceaux
séparément en chaque variable — ce que la page appelle des
\emph{bifaisceaux}.\footnote{La page écrit la troisième condition avec un mot
biffé et non remplacé, et la quatrième description du côté $D$ avec un membre
illisible ; les deux se rétablissent sans ambiguïté par symétrie avec les deux
premières, que la page écrit en toutes lettres. C'est aussi la seule lecture
qui rende vraie la Proposition 1 ci-dessous.}
\end{itemize}

La dernière description est celle qui porte tout : un objet du produit est une
donnée locale à deux variables, locale séparément en chacune.

\pagerange{3}{4}
\subsection*{$\Pi(E,F)$ est un topos}

\textbf{Proposition 1.} \emph{$\Pi(E,F)$ est un topos.}

La démonstration est un emboîtement de sous-topos, et elle est complète. On
peut supposer $E = \widetilde{C}$, $F = \widetilde{D}$ avec $C$, $D$ petites et
stables par limites projectives finies.

\begin{enumerate}
\item[a)] Topologies discrètes. Un bifaisceau sur $(\widehat{C}, \widehat{D})$
n'est astreint à rien, donc
\[
\Pi(\widehat{C}, \widehat{D}) \;\simeq\; \widehat{C \times D},
\]
qui est un topos.
\item[b)] Un cran plus loin, $\Pi(\widehat{C}, \widetilde{D}) \simeq
\operatorname{Hom}(C^{\circ}, \widetilde{D})$, et l'inclusion
\[
\Pi(\widehat{C}, \widetilde{D}) \hookrightarrow
\Pi(\widehat{C}, \widehat{D}), \qquad \varphi \longmapsto i_{D} \circ \varphi,
\]
est le foncteur image directe d'un plongement de topos : son adjoint à gauche
est $\varphi \mapsto a_{D} \circ \varphi$, qui est exact à gauche puisque
$a_{D}$ l'est et que limites finies et composition se calculent argument par
argument. Donc $\Pi(\widehat{C}, \widetilde{D})$ est un \emph{sous-topos} de
$\Pi(\widehat{C}, \widehat{D})$, et $\Pi(\widetilde{C}, \widehat{D})$ aussi,
par la raison symétrique.
\end{enumerate}

Enfin
\[
\Pi(\widetilde{C}, \widetilde{D}) \;=\; \Pi(\widetilde{C}, \widehat{D})
\,\cap\, \Pi(\widehat{C}, \widetilde{D})
\]
à l'intérieur de $\Pi(\widehat{C}, \widehat{D})$ : un bifaisceau est un
faisceau en chaque variable, ce qui est exactement l'appartenance aux deux
sous-topos à la fois. Or l'intersection de deux sous-topos est un sous-topos.
Donc $\Pi(E,F)$ est un topos.\footnote{Ce dernier pas est le seul que la page
invoque sans le démontrer, et c'est aujourd'hui un énoncé standard :
les sous-topos d'un topos $\mathcal{E}$ correspondent bijectivement aux
topologies de Lawvere--Tierney sur $\mathcal{E}$, et forment un treillis
complet, l'intersection correspondant à la borne supérieure des topologies.
Le manuscrit note l'inclusion « str. pleine » comme une lecture douteuse ;
c'est bien de sous-catégories strictement pleines qu'il s'agit.}

\pagerange{4}{5}
\subsection*{Fonctorialité}

Pour deux morphismes de topos $u : E \to E'$ et $v : F \to F'$, on veut
$\Pi(u,v) : \Pi(E,F) \to \Pi(E',F')$. En termes de bifaisceaux, l'image directe
est immédiate :
\[
\Pi(u,v)_{*}(\varphi)(X',Y') \;=\; \varphi(u^{*}X',\, v^{*}Y').
\]
La page laisse en regard « $\Pi(u,v)^{*}(\varphi)(X,Y) \overset{?}{=}$ » sans
la remplir.\footnote{Elle est déterminée par adjonction et n'a pas d'expression
aussi simple. Sous l'équivalence $\Pi(E,F) \simeq E \times F$ établie plus bas,
$\Pi(u,v)$ n'est autre que $u \times v$, et son image inverse est le produit
des images inverses — ce qui est la réponse au point d'interrogation, mais elle
suppose acquis ce que la page n'a pas encore démontré.}

Que ce couple de foncteurs soit bien un morphisme de topos se ramène, en
factorisant $\Pi(E,F) \to \Pi(E,F') \to \Pi(E',F')$, au cas où l'un des deux
morphismes est l'identité. Restent alors deux vérifications, dont la seconde
est le cœur : sur $\operatorname{Hom}(C, F) \to \operatorname{Hom}(C, F')$, le
foncteur en jeu est $\varphi \mapsto v_{*} \circ \varphi$, et son \emph{adjoint
à gauche} est $\varphi \mapsto v^{*} \circ \varphi$, qui est exact à gauche
puisque $v^{*}$ l'est. C'est bien un morphisme de topos.\footnote{La page écrit
« dont l'adjoint à droite est $\varphi \mapsto v^{*} \circ \varphi$ ». C'est un
lapsus : $v^{*} \dashv v_{*}$, donc $v^{*} \circ -$ est adjoint \emph{à gauche}
de $v_{*} \circ -$, et c'est bien l'adjoint à gauche qu'il faut ici, puisque
c'est de lui qu'on demande l'exactitude à gauche. La transcription le signale ;
l'argument de la page est juste, seul le mot est faux.}

\pagerange{5}{8}
\subsection*{C'est le 2-produit}

Prenant pour $v$ le morphisme $F \to P$ vers le topos ponctuel et notant que
$\Pi(E,P) \simeq E$, on obtient deux projections canoniques $\Pi(E,F) \to E$ et
$\Pi(E,F) \to F$, d'où un morphisme de topos
\[
\Pi(E,F) \longrightarrow E \times_{\mathrm{top}} F,
\]
2-fonctoriel en $E$ et $F$. \emph{C'est une équivalence}, et c'est là le
théorème : $\Pi(E,F)$ « est » le 2-produit des topos $E$ et $F$.

La démonstration se réduit encore. On a vu que
$\Pi(\widetilde{C}, \widetilde{D})$ est l'intersection, dans
$\Pi(\widehat{C}, \widehat{D}) \simeq \widehat{C \times D}$, des images
inverses des sous-topos $\widetilde{C} \subset \widehat{C}$ et
$\widetilde{D} \subset \widehat{D}$ par les deux projections. Il suffit donc de
savoir que $\widehat{C \times D}$ est le 2-produit de $\widehat{C}$ et
$\widehat{D}$, et cela résulte du calcul suivant : pour tout topos $X$,
\[
\operatorname{Hom}_{\mathrm{top}}(X, \widehat{C}) \;\simeq\;
\operatorname{Hom}_{\mathrm{ex.g.}}(C, X)^{\circ},
\]
où $\operatorname{Hom}_{\mathrm{ex.g.}}$ désigne les foncteurs exacts à
gauche.\footnote{Deux précisions que la page porte en marge, et qui sont
nécessaires. D'abord l'hypothèse : $C$ doit être stable par limites projectives
finies, faute de quoi le membre de droite est la catégorie des foncteurs plats
et non celle des foncteurs exacts à gauche. Ensuite l'exposant $(\ )^{\circ}$,
qui enregistre une convention et non un fait : selon que l'on prend pour
2-flèches de topos les transformations des images directes ou celles des images
inverses, on obtient l'une ou l'autre des deux catégories opposées. Le choix ne
change ni les objets ni l'énoncé du 2-produit.} On a alors
\[
\operatorname{Hom}_{\mathrm{ex.g.}}(C, X) \times
\operatorname{Hom}_{\mathrm{ex.g.}}(D, X) \;\simeq\;
\operatorname{Hom}_{\mathrm{ex.g.}}(C \times D, X),
\]
qui est le point de fond : un foncteur exact à gauche $w$ sur $C \times D$
équivaut au couple $\bigl(a \mapsto w(a, e_{D}),\ b \mapsto w(e_{C}, b)\bigr)$,
de reconstitution $w(a,b) \simeq u(a) \times v(b)$. La vérification que ce $w$
est encore exact à gauche est le calcul que la page mène en marge, et il n'y a
là de subtil qu'un point : les limites dans $C \times D$ se calculent
composante par composante, mais elles portent sur des diagrammes de
\emph{couples} et non sur les deux variables séparément — mise en garde que la
page prend soin d'écrire.

Deux corollaires suivent :

\begin{enumerate}
\item[a)] pour une famille $(C_{i})_{i \in I}$ de petites catégories à limites
finies,
\[
\prod_{i} \widehat{C_{i}} \;\simeq\; \widehat{\textstyle\prod'_{i} C_{i}},
\]
où $\prod'$ est le \emph{produit restreint} : la sous-catégorie pleine du
produit formée des systèmes $(x_{i})$ tels que $x_{i} \simeq e_{C_{i}}$ pour
presque tout $i$ ;
\item[b)] si $E'_{i} \subset E_{i}$ sont des sous-topos et $E = \prod_{i} E_{i}$,
alors $\bigcap_{i} \mathrm{pr}_{i}^{-1}(E'_{i})$ est un 2-produit des $E'_{i}$.
\end{enumerate}

\textbf{Corollaire.} \emph{Le 2-produit d'une famille quelconque de topos
existe toujours.}\footnote{Énoncé aujourd'hui classique — la 2-catégorie des
topos de Grothendieck admet toutes les petites 2-limites — et dont c'est ici la
démonstration par les sites. La page 8 entreprend ensuite le cas général d'une
$\varprojlim$ de topos fibré, écrit le diagramme des trois étages
$\prod_{i} E_{i} \rightrightarrows \prod_{\alpha} E_{b(\alpha)}
\rightrightarrows \prod_{(\alpha,\beta)} E_{b(\beta)}$, puis barre tout le bas
de la page de deux longues diagonales. Nous ne restituons pas un argument
qu'il a lui-même abandonné ; la page 15 y reviendra sous une hypothèse
filtrante.}

\pagerange{9}{9}
\section*{Le produit d'espaces contre le produit de topos (pages 9 à 12)}

Soient $X$ et $Y$ deux espaces topologiques, $\mathcal{O}_{X}$ la catégorie
(l'ensemble ordonné) de leurs ouverts. Le produit cartésien $X \times Y$ porte
sa topologie usuelle, dont les rectangles $U \times V$ forment une base ; d'où
un morphisme canonique de topos
\[
(*) \qquad \operatorname{Top}(X \times Y) \longrightarrow
\operatorname{Top}(X) \times_{\mathrm{top}} \operatorname{Top}(Y),
\]
et la question : \emph{quand est-ce une équivalence ?}

La réponse ne va pas de soi, et c'est ce qui rend ces quatre pages
intéressantes. Le membre de droite est, par ce qui précède, le topos des
bifaisceaux sur $\mathcal{O}_{X} \times \mathcal{O}_{Y}$ ; le membre de gauche
est le topos des faisceaux sur $X \times Y$. Les deux sont des topos de
faisceaux sur le \emph{même} site $\mathcal{O}_{X} \times \mathcal{O}_{Y}$,
mais pour deux topologies distinctes :\footnote{Cette identification suppose
$X \neq \emptyset$ et $Y \neq \emptyset$, pour que
$\mathcal{O}_{X} \times \mathcal{O}_{Y} \subset \mathcal{O}_{X \times Y}$ soit
une injection ; la page le suppose expressément. Elle note aussi, sans le
souligner, que $\operatorname{Top}(X \times Y)$ ne dépend que de
$\operatorname{Top}(X)$ et $\operatorname{Top}(Y)$, c'est-à-dire des espaces
sobres associés — fait que les pages 34 et 36 reprendront pour lui-même.}

\begin{itemize}
\item la topologie \emph{produit} $\pi$, dont les faisceaux sont les
bifaisceaux ;
\item la topologie $\pi'$ induite par celle de $X \times Y$, où recouvrir veut
dire recouvrir ensemblistement.
\end{itemize}

Tout faisceau sur $X \times Y$ est un bifaisceau, donc $\pi'$ est plus fine que
$\pi$, et $(*)$ est un \emph{plongement} de topos. Reste à savoir s'il est une
équivalence, c'est-à-dire si $\pi = \pi'$, c'est-à-dire enfin si tout
recouvrement ensembliste
\[
\bigcup_{i} U_{i} \times V_{i} \;=\; U \times V
\qquad\text{entraîne}\qquad
\operatorname*{Sup}_{i} U_{i} \times V_{i} \;=\; U \times V,
\]
le $\operatorname{Sup}$ étant pris dans le treillis des sous-objets de
$U \times V$ dans le topos produit.\footnote{C'est exactement l'écart, familier
en algèbre linéaire, entre les fonctions de deux variables et celles qu'on
obtient à partir des produits $f(x)g(y)$ : les rectangles engendrent la
topologie de $X \times Y$ comme base, mais ce qu'ils engendrent par
\emph{sup dans le produit de topos} peut être strictement plus petit. En
langage d'aujourd'hui, $(*)$ compare le produit d'espaces au produit de
\emph{locales} $\mathcal{O}_{X} \otimes \mathcal{O}_{Y}$, et la question est
celle de la spatialité du second.}

\pagerange{9}{10}
\subsection*{Le critère, et sa démonstration}

\textbf{Proposition.} \emph{Pour qu'il en soit ainsi, il suffit que tout point
de $U$ — ou tout point de $V$ — admette un voisinage quasi-compact contenu dans
l'ouvert.}

\textbf{Corollaire.} \emph{Si tout point de $X$ (ou de $Y$) admet un système
fondamental de voisinages quasi-compacts, $(*)$ est une équivalence.}

La démonstration tient en trois passages à la borne supérieure, et il vaut de
la donner : elle est complète sur la page, et c'est le lemme du tube.

Pour $x \in X$, soit $U_{x} \subset U$ un voisinage quasi-compact de $x$, et
$\mathring{U}_{x}$ son intérieur. Fixons $y \in V$. Comme $U_{x} \times \{y\}$
est quasi-compact, il est recouvert par un nombre \emph{fini} des
$U_{i} \times V_{i}$, disons pour $i \in I_{x,y}$ ; posons
$V_{x,y} = \bigcap_{i \in I_{x,y}} V_{i}$, voisinage ouvert de $y$. On a alors
$U_{x} \subset \bigcup_{i \in I_{x,y}} U_{i}$ et $V_{i} \supset V_{x,y}$ pour
$i \in I_{x,y}$, d'où, en posant
$S = \operatorname*{Sup}_{i \in I}(U_{i} \times V_{i})$ :
\[
S \;\geqslant\; \operatorname*{Sup}_{i \in I_{x,y}}(U_{i} \times V_{i})
\;\geqslant\; \operatorname*{Sup}_{i \in I_{x,y}}(U_{i} \times V_{x,y})
\;=\; \Bigl(\operatorname*{Sup}_{i \in I_{x,y}} U_{i}\Bigr) \times V_{x,y}
\;\geqslant\; \mathring{U}_{x} \times V_{x,y}.
\]
La seule égalité de cette chaîne est celle qui fait marcher tout l'argument :
$W \mapsto W \times V'$ commute aux bornes supérieures quelconques dans le
produit de locales.\footnote{C'est vrai, et c'est précisément \emph{ce que} le
produit de topos garantit : les sous-objets de l'objet final y forment un
treillis où les $\inf$ finis sont distributifs par rapport aux $\sup$
quelconques. La page ne le justifie pas ; la page 28 y revient en propres
termes.} Faisant maintenant varier $y$, puis $x$ :
\[
S \;\geqslant\; \operatorname*{Sup}_{y \in V}
\bigl(\mathring{U}_{x} \times V_{x,y}\bigr)
= \mathring{U}_{x} \times V,
\qquad
S \;\geqslant\; \operatorname*{Sup}_{x \in U}
\bigl(\mathring{U}_{x} \times V\bigr) = U \times V,
\]
puisque les $V_{x,y}$ recouvrent $V$ et les $\mathring{U}_{x}$ recouvrent $U$.
D'où l'égalité cherchée.

\pagerange{10}{11}
\subsection*{Le sens réciproque}

La page 10 poursuit par une réciproque, sous une hypothèse plus faible que la
compacité locale : \emph{supposons que toute partie localement fermée non vide
de $X$ possède un point admettant, dans elle, un voisinage relatif
quasi-compact.} Alors $(*)$ est encore une équivalence.

L'argument est un argument de maximalité, et il faut le compléter d'un pas que
la page passe. Soit $U'$ le plus grand ouvert de $U$ tel que
$S \geqslant U' \times V$ — il existe, la famille de ces ouverts étant stable
par réunion et contenant $\emptyset$. Supposons $U' \neq U$. Alors
$Z = U \setminus U'$ est localement fermé et non vide, donc contient un point
$x$ admettant dans $Z$ un voisinage relatif quasi-compact $T$ ; soit
$U'' \subset U$ un ouvert contenant $x$ avec $U'' \cap Z \subset T$. Pour
chaque $y \in V$, la quasi-compacité de $T$ donne une partie finie $I_{y}$ avec
$T \subset \bigcup_{i \in I_{y}} U_{i}$, et $V_{y} = \bigcap_{i \in I_{y}} V_{i}$
est un voisinage de $y$ ; comme au paragraphe précédent,
\[
S \;\geqslant\; \Bigl(\operatorname*{Sup}_{i \in I_{y}} U_{i}\Bigr) \times V_{y}.
\]
Or $U'' \setminus U' \subset U'' \cap Z \subset T \subset
\bigcup_{i \in I_{y}} U_{i}$, donc $U'' \subset U' \cup
\bigcup_{i \in I_{y}} U_{i}$, et comme par ailleurs
$S \geqslant U' \times V \geqslant U' \times V_{y}$, il vient
$S \geqslant (U' \cup U'') \times V_{y}$.\footnote{Cette inclusion
$U'' \subset U' \cup \bigcup_{i \in I_{y}} U_{i}$ est le pas que la page ne
formule pas : elle écrit directement
$S \geqslant \operatorname{Sup}(U', U_{i}\ i \in I_{y}) \times V_{y}
\geqslant (U' \cup U'') \times V_{y}$. Elle est immédiate, et sans elle la
conclusion ne suit pas. Le reste de la page est correct tel quel.} Faisant
varier $y$, on obtient $S \geqslant (U' \cup U'') \times V$, ce qui contredit
la maximalité de $U'$ puisque $x \in U'' \setminus U'$. Donc $U' = U$.

\pagerange{11}{14}
\subsection*{Produits infinis}

Les pages 11 et 12 passent à une famille quelconque $(X_{i})_{i \in I}$, avec
la conclusion suivante : si tous les facteurs sauf au plus un ont des systèmes
fondamentaux de voisinages quasi-compacts, et si tous sauf au plus un sont
quasi-compacts, alors
\[
\operatorname{Top}\Bigl(\prod_{i \in I} X_{i}\Bigr) \longrightarrow
\prod_{i}{}^{\mathrm{top}} \operatorname{Top}(X_{i})
\]
est une équivalence.\footnote{C'est aujourd'hui l'énoncé standard : un produit
de locales est spatial et coïncide avec le produit d'espaces dès que presque
tous les facteurs sont compacts et les autres localement compacts. La double
hypothèse du manuscrit — « tous sauf au plus un » deux fois — est exactement
celle qu'il faut pour que l'argument du tube passe à la limite, et elle n'est
pas plus forte que nécessaire.}

L'argument est le même, mené sur le \emph{produit restreint}
$\prod' \mathcal{O}_{X_{i}}$ — les familles $(U_{i})$ avec $U_{i} = X_{i}$ pour
presque tout $i$ — et sur $I = \varinjlim_{\alpha} J_{\alpha}$, réunion filtrante
de ses parties finies. Le cas fini étant acquis, on écrit
\[
\prod_{i}{}^{\mathrm{top}} \operatorname{Top}(X_{i})
\;\simeq\; \varprojlim_{\alpha}{}^{\mathrm{top}}
\operatorname{Top}\bigl(\textstyle\prod_{i \in J_{\alpha}} X_{i}\bigr)
\;\simeq\; \bigl(\varinjlim_{\alpha} \mathcal{O}_{Z_{\alpha}}\bigr)^{\sim},
\]
et l'on reprend le même dévissage : tout $W$ du produit restreint est de la
forme $U \times Z'_{\alpha}$ avec $Z'_{\alpha} = \prod_{i \notin J_{\alpha}} X_{i}$
quasi-compact, et une famille $(V_{\lambda})$ le recouvrant se raffine, en
chaque point, sur un $J_{\beta} \supset J_{\alpha}$ fini.

La page 14 aborde enfin les \emph{contre-exemples}, avec tous les $X_{i}$
discrets, écrit le système fondamental de voisinages
$V_{z} = \{z' \mid z'_{i} = z_{i}\ \text{hors d'une partie finie}\}$, se demande
si ces $V_{z}$ recouvrent pour $\pi$ — et s'arrête, un point d'interrogation en
marge.\footnote{La question est réelle et la réponse négative en général : sans
compacité, le produit de locales cesse d'être spatial, et le morphisme $(*)$
est un plongement strict. C'est le phénomène qu'Isbell a isolé et qui a fait
des locales un objet d'étude pour lui-même. Nous n'attribuons pas au manuscrit
un contre-exemple qu'il n'écrit pas : il s'arrête sur son point
d'interrogation, et nous avec lui.}

\pagerange{15}{15}
\section*{Limites filtrantes de topos essentiels (page 15)}

Une note d'une page, indépendante de ce qui précède, et laissée inachevée.

Un morphisme de topos est dit \emph{essentiel} lorsque son image inverse admet
elle-même un adjoint à gauche. Pour les topos de préfaisceaux, ce sont
exactement les morphismes provenant d'un foncteur entre les sites. Soit donc un
système projectif filtrant de topos $\widehat{C_{i}}$, dont les morphismes de
transition $\widehat{f}_{ij} : \widehat{C_{j}} \to \widehat{C_{i}}$ proviennent
de foncteurs $f_{ij} : C_{j} \to C_{i}$ admettant des adjoints à droite
$g_{ji} : C_{i} \to C_{j}$. Alors
\[
\mathcal{X} \;=\; \varprojlim_{i}{}^{\mathrm{top}} \widehat{C_{i}}
\;\simeq\; \widehat{C}, \qquad C = \operatorname*{colim}_{i, g} C_{i},
\]
et les projections $\mathrm{pr}_{i} : \mathcal{X} \to \widehat{C_{i}}$ sont
associées aux foncteurs canoniques $\alpha_{i} : C_{i} \to C$, avec
\[
\mathrm{pr}_{i*}(F) = F \circ \alpha_{i}, \qquad
\mathrm{pr}_{i}^{*}(G) = \alpha_{i!}(G),
\]
$\alpha_{i!}$ désignant l'extension de Kan à gauche le long de
$\alpha_{i}$.\footnote{La page s'arrête sur « $\mathrm{pr}_{i}^{*}(G) =$ » :
la formule n'est pas écrite. Celle donnée ici est la nôtre — elle est
déterminée sans ambiguïté par l'adjonction, l'image directe étant la
restriction le long de $\alpha_{i}$. Nous la donnons parce qu'elle est forcée,
et non parce que la page la suggère.} \emph{Une limite filtrante de topos se
calcule comme une colimite filtrante de sites} : c'est l'énoncé, et l'hypothèse
d'essentialité est ce qui permet de le lire au niveau des sites plutôt qu'à
celui des topos.

\pagerange{17}{17}
\section*{Les foncteurs cocontinus sur un topos de faisceaux (pages 17 à 20)}

La dernière note de cette série est la plus achevée, et c'est elle qui court
jusqu'à la page 21. Soient $C$ un site et $E$ une catégorie où les petites colimites
existent. On veut décrire les foncteurs $u : \widetilde{C} \to E$ commutant aux
petites colimites, à partir de leur seule restriction à $C$.

\pagerange{17}{18}
\subsection*{L'adjonction}

Notons $\varepsilon_{C} = a_{C} \circ j_{0} : C \to \widetilde{C}$ le foncteur
canonique ($j_{0} : C \to \widehat{C}$ étant Yoneda), et posons
\[
\alpha(u) = u \circ \varepsilon_{C}, \qquad
\beta(u_{0}) = \Bigl[\, F \longmapsto \varinjlim_{C/F} u_{0}(X) \,\Bigr],
\]
deux foncteurs entre $\operatorname{Hom}(\widetilde{C}, E)$ et
$\operatorname{Hom}(C, E)$.\footnote{Les pages 19 et 20 notent $a$ le foncteur
$u \mapsto u \circ \varepsilon_{C}$, entrant en collision avec le $a$ du
faisceau associé, qui est employé sur la même page. Nous gardons $\alpha$,
$\beta$ des pages 17 et 18, et réservons $a$ à la sheafification. La page 19
écrit en outre $\beta : \operatorname{Hom}(C,E) \to \operatorname{Hom}(C,E)$,
le but devant être $\operatorname{Hom}(\widetilde{C},E)$.}

Deux morphismes fonctoriels s'écrivent alors sans effort. Le premier,
\[
i_{u} : \beta\alpha(u) \longrightarrow u,
\qquad
\varinjlim_{C/F} u(\varepsilon_{C}X) \longrightarrow u(F),
\]
est déduit de l'isomorphisme canonique
$\varinjlim_{C/F} \varepsilon_{C}(X) \xrightarrow{\ \sim\ } F$ en lui
appliquant $u$. Le second,
\[
j_{u_{0}} : u_{0} \longrightarrow \alpha\beta(u_{0}),
\]
envoie $u_{0}(Z)$ dans $\varinjlim_{C/\varepsilon_{C}(Z)} u_{0}(X)$ par la
composante d'indice $Z$ lui-même. Ce sont les morphismes d'adjonction d'une
adjonction $\beta \dashv \alpha$ : $i$ en est la coünité, $j$ l'unité.

Et l'on lit aussitôt sur eux la partie fixe de l'adjonction, au sens usuel :

\begin{itemize}
\item $i_{u}$ est un isomorphisme si et seulement si $u$ commute aux colimites
de la forme $\varinjlim_{C/F} X \xrightarrow{\sim} F$, c'est-à-dire si et
seulement si $u$ est cocontinu ;
\item $j_{u_{0}}$ est un isomorphisme si et seulement si, pour tout
$Z \in C$, $u_{0}(Z) \to \varinjlim_{C/\varepsilon_{C}(Z)} u_{0}(X)$ l'est —
condition \emph{vide} dès que $\varepsilon_{C}$ est pleinement fidèle,
c'est-à-dire dès que la topologie de $C$ est moins fine que la topologie
canonique, car $C/\varepsilon_{C}(Z) \simeq C/Z$ admet alors $Z$ pour objet
final.
\end{itemize}

L'adjonction se restreint donc en une équivalence entre ses deux parties fixes,
et en particulier $\alpha$ induit un foncteur pleinement fidèle
$\operatorname{Hom}'(\widetilde{C}, E) \to \operatorname{Hom}(C, E)$, où
$\operatorname{Hom}'$ désigne partout les foncteurs commutant aux petites
colimites.

\pagerange{19}{20}
\subsection*{L'image essentielle, et la question du bas de page}

Reste à décrire cette image essentielle, et c'est ici que les pages 19 et 20
prennent le détour qui porte. On identifie d'abord
$\operatorname{Hom}(C, E) \simeq \operatorname{Hom}'(\widehat{C}, E)$ —
propriété universelle de $\widehat{C}$ comme cocomplétion libre de $C$ — ce qui
transporte $\alpha$ et $\beta$ en
\[
\alpha'(u) = u \circ a, \qquad \beta'(v) = v \mid \widetilde{C},
\]
avec $\beta'\alpha'(u) \simeq u$. La pleine fidélité de $\alpha'$ tient à ce
que $\widetilde{C}$ est, par $a$, une catégorie de fractions de
$\widehat{C}$ : $\operatorname{Hom}(\widetilde{C}, E)$ s'identifie à la
sous-catégorie pleine de $\operatorname{Hom}(\widehat{C}, E)$ des foncteurs qui
inversent les flèches que $a$ inverse — les morphismes \emph{bicouvrants}. Et
la cocontinuité se transporte dans cette identification, par le calcul de la
page 20 : si $\mathcal{U} \circ a$ commute aux colimites, alors
\[
\mathcal{U}\Bigl(\varinjlim_{\alpha}{}^{\widetilde{C}} X_{\alpha}\Bigr)
= \mathcal{U}a\Bigl(\varinjlim_{\alpha}{}^{\widehat{C}} X_{\alpha}\Bigr)
= \varinjlim_{\alpha} (\mathcal{U}a)(X_{\alpha})
= \varinjlim_{\alpha} \mathcal{U}(X_{\alpha}),
\]
puisque les colimites de $\widetilde{C}$ sont celles de $\widehat{C}$ suivies
de $a$.

D'où l'énoncé, qui est la conclusion de la note :

\textbf{Théorème.} \emph{$\alpha'$ induit une équivalence de
$\operatorname{Hom}'(\widetilde{C}, E)$ sur la sous-catégorie pleine de
$\operatorname{Hom}(\widehat{C}, E)$ formée des foncteurs qui} a) \emph{commutent
aux petites colimites, et} b) \emph{transforment les morphismes bicouvrants en
isomorphismes.}\footnote{C'est la propriété universelle du topos de faisceaux
comme localisation exacte à gauche de $\widehat{C}$ : $\widetilde{C}$ est la
cocomplétion libre de $C$ \emph{modulo} les colimites que la topologie
impose. Sous cette forme, l'énoncé est aujourd'hui celui qu'on donne des
foncteurs « continus » d'un site vers une catégorie cocomplète.}

La page 20 s'achève sur une tentative d'affaiblissement. La condition b) est
malcommode : elle porte sur $\widehat{C}$ et non sur $C$. Ne suffirait-il pas
de demander que toute famille couvrante $X_{\alpha} \to X$ aille sur une
famille épimorphique de $E$ ? « \emph{Cela est-il suffisant ?} », écrit-il —
et, en bas de page, au crayon, \emph{non}.

Le « non » est juste. Une famille épimorphique ne contrôle que la surjectivité
du morphisme de comparaison $\varinjlim_{C/R} u_{0} \to u_{0}(X)$, alors que
b) en demande la bijectivité : un crible couvrant $R \subset X$ doit aller sur
un \emph{isomorphisme}, ce qui exige en outre le recollement, c'est-à-dire
l'injectivité.\footnote{La page ne donne pas de contre-exemple, et nous n'en
fabriquons pas : ce qui est ici affirmé est que la condition proposée est
strictement plus faible que b), ce qui se lit sur les deux énoncés. La page 21
porte la phrase suivante — « soit $F \to G$ dans $\widehat{C}$,
bicouvrant, prouvons que \ldots\ On a » — et s'interrompt immédiatement, la
ligne d'après étant entièrement raturée. La tentative a été abandonnée là.}

\pagerange{21}{21}
\section*{La phrase interrompue (page 21)}

La page 20 s'achevait sur une question — la condition « toute famille
couvrante va sur une famille épimorphique » suffit-elle à caractériser les
foncteurs cocontinus sur $\widetilde{C}$ ? — et sur un « non » porté au crayon
en bas de page. La page 21 en commence la démonstration : « soit
$F \to G$ dans $\widehat{C}$, bicouvrant, prouvons que
$\widehat{u}_{0}(F) \to \widehat{u}_{0}(G)$ est un isomorphisme. On a \ldots »
La ligne suivante est entièrement raturée, et rien ne poursuit.

Nous n'y suppléons pas. La raison pour laquelle la réponse est négative a été
donnée à sa place, plus haut : une famille épimorphique ne contrôle que la
surjectivité du morphisme de comparaison, là où le caractère bicouvrant en
demande la bijectivité.


\pagerange{23}{23}
\section*{Espaces essentiels et ensembles préordonnés (pages 23 à 26)}

\pagerange{23}{23}
\subsection*{Le dictionnaire}

Soit $X$ un ensemble, $\mathcal{T}(X)$ l'ensemble des topologies sur $X$ et
$\underline{\omega}(X)$ celui des préordres. Deux constructions les relient,
fonctorielles en $X$ pour les images inverses.

\emph{D'une topologie à un préordre.} On pose
\[
x \leq y \iff x \in \overline{\{y\}}
\]
— le \emph{préordre de spécialisation} : $x$ est une spécialisation de $y$, et
$y$ le point le plus générique des deux.\footnote{C'est ici qu'une convention
doit être fixée pour tout le dossier, et le manuscrit ne la tient pas. La page
23 définit $\alpha_{X}(T)$ par « $y \leq x \iff x \in \overline{\{y\}}$ »,
c'est-à-dire l'inverse de ce qui précède ; mais les trois autres définitions de
la même page — les ouverts comme parties stables vers le haut, les fermés vers
le bas, et $O_{x} = \{y \mid y \geq x\}$ — s'accordent entre elles et exigent
la convention adoptée ici. Lue littéralement, la page donne
$\alpha\beta(\omega) = \omega^{\circ}$ et contredit son propre énoncé
$\alpha\beta(\omega) = \omega$. Nous échangeons donc les deux lettres dans la
définition de $\alpha$, et tout le reste de la page devient exact sans autre
retouche. Toutes les inégalités de cette section sont écrites dans la
convention corrigée, y compris là où la page écrit la relation opposée.}

\emph{D'un préordre à une topologie.} On déclare ouvertes les parties stables
vers le haut,
\[
\mathrm{Ouv} = \{\, U \subset X \mid x \in U,\ x \leq y \Rightarrow y \in U \,\},
\]
qui sont exactement les réunions quelconques des
$O_{x} = \{\, y \mid y \geq x \,\}$ ; les fermés sont les parties stables vers
le bas. C'est ce qu'on appelle aujourd'hui la \emph{topologie d'Alexandrov} du
préordre, et $O_{x}$ y est le plus petit voisinage ouvert de $x$.

Une application entre espaces topologiques qui est continue est croissante ;
une application entre préordres qui est croissante est continue pour les
topologies associées. D'où deux foncteurs, et deux morphismes fonctoriels :

\begin{itemize}
\item $\rho : \beta\alpha \to \mathrm{id}_{(\mathrm{Esp})}$, qui traduit le fait
que \emph{tout ouvert est stable par générisation} — c'est-à-dire que la
topologie de départ est moins fine que la topologie d'Alexandrov de son
préordre de spécialisation ;
\item $\sigma : \mathrm{id}_{(\text{Préord})} \xrightarrow{\ \sim\ } \alpha\beta$,
qui est un \emph{isomorphisme} : le préordre de spécialisation de la topologie
d'Alexandrov de $\omega$ est $\omega$.
\end{itemize}

Ce sont les morphismes d'adjonction d'une adjonction $\beta \dashv \alpha$ :
\[
\operatorname{Hom}_{(\mathrm{Esp})}(\beta I, X) \;\simeq\;
\operatorname{Hom}_{(\text{Préord})}(I, \alpha X).
\]

\pagerange{24}{24}
\subsection*{Ce que le dictionnaire attrape}

$\sigma$ étant un isomorphisme, $\beta$ est \emph{pleinement fidèle} : les
préordres forment une sous-catégorie \emph{coréflexive} de celle des espaces.
Son image essentielle est décrite exactement :

\textbf{Corollaire.} \emph{$\beta$ est pleinement fidèle, et son image
essentielle est formée des espaces où tout point $x$ possède un plus petit
voisinage ouvert $O_{x}$ — de façon équivalente, où toute intersection
d'ouverts est ouverte.}

La seule chose à vérifier est que $O_{x}$, lorsqu'il existe, coïncide avec
$\{y \mid y \geq x\}$, et cela tient en une ligne : $y \in O_{x}$ signifie que
tout ouvert contenant $x$ contient $y$, c'est-à-dire $x \in \overline{\{y\}}$.
Ces espaces sont ceux qu'on appelle aujourd'hui \emph{d'Alexandrov} ; le
manuscrit les dira \emph{essentiels}, pour une raison qui apparaîtra plus
bas.\footnote{La page note en marge que cette classe n'est pas stable par une
opération dont le nom est illisible. C'est le produit : le produit d'une
famille infinie d'espaces d'Alexandrov n'est pas d'Alexandrov, une
intersection infinie d'ouverts du produit cessant d'être ouverte. La page 24
signale d'ailleurs, dans une note portée en diagonale, que $\beta$ est exact à
gauche — la topologie associée à un préordre produit \emph{fini} est bien le
produit des topologies — mais ne commute pas aux produits quelconques, et
annonce un contre-exemple qu'elle n'écrit pas.}

Deux compléments de la même page :

\begin{itemize}
\item la relation de spécialisation est une relation d'\emph{ordre} — c'est-à-dire
antisymétrique — si et seulement si $X$ vérifie $\overline{\{x\}} =
\overline{\{y\}} \Rightarrow x = y$ ;\footnote{La page appelle cette condition
la \emph{sobriété}, et la transcription signale le mot comme une lecture
douteuse. Sous cette lecture l'énoncé est faux : la condition écrite est
l'axiome $T_{0}$, strictement plus faible que la sobriété, qui exige en outre
que tout fermé irréductible ait un point générique. Le manuscrit lui-même est
exact sur ce point ailleurs : la page 34 distingue explicitement les deux
moitiés — injectivité pour $T_{0}$, surjectivité pour l'existence des points
génériques, bijectivité pour la sobriété.}
\item pour $X$ d'Alexandrov, en notant $\underline{X}$ la catégorie du
préordre (une flèche $x \to y$ lorsque $x \leq y$), le foncteur
$F \mapsto (x \mapsto F(O_{x}))$ est une \emph{équivalence}
\[
\operatorname{Top}(X) \;\simeq\; \operatorname{Hom}(\underline{X},
(\mathrm{Ens})),
\]
la catégorie des foncteurs \emph{covariants} sur $\underline{X}$.\footnote{La
page écrit $\widehat{\underline{X}}$, c'est-à-dire les foncteurs
contravariants. La différence est un passage à l'opposé, et elle est forcée par
la convention choisie plus haut : si $x \leq y$ alors $O_{y} \subset O_{x}$,
donc $x \mapsto F(O_{x})$ est covariant en $x$. Avec la convention opposée
c'est la forme du manuscrit qui serait juste ; l'une ou l'autre convient,
pourvu qu'on n'en change pas en route.} Sous cette équivalence, les ouverts de
$X$ correspondent aux sous-objets de l'objet final, ce qui est la description
d'un \emph{topos localique}.
\end{itemize}

\pagerange{25}{26}
\subsection*{Points, points génériques, points essentiels}

La page 25 poursuit sur les fermés irréductibles. Pour $X$ d'Alexandrov, $X$
est irréductible si et seulement si deux ouverts non vides se rencontrent,
c'est-à-dire si et seulement si l'ensemble ordonné est \emph{filtrant}, et les
fermés irréductibles sont donc les parties $Y$ qui sont a) stables vers le bas
et b) filtrantes.\footnote{La page écrit « cofiltrant » et « $\exists z$ avec
$z \leq x$, $z \leq y$ », ce qui est la même condition lue dans la convention
opposée de la note précédente. Nous écrivons filtrant, dans la nôtre :
$O_{x} \cap O_{y} \neq \emptyset$ équivaut à l'existence d'un $z \geq x$,
$z \geq y$.} On retrouve ainsi le \emph{sobrifié} $\widetilde{X}$ — l'ensemble
des fermés irréductibles, ordonné — et l'application canonique
\[
x \longmapsto \overline{\{x\}} = \{\, y \mid y \leq x \,\},
\qquad X \longrightarrow \widetilde{X},
\]
qui s'identifie au foncteur canonique $\underline{X} \to
\mathrm{Pro}(\underline{X})$ et à l'application de $X$ dans l'espace sobre
associé.

Un point générique d'un fermé irréductible $Y$ est alors un \emph{plus grand
élément} de $Y$, et la page 26 énumère les conditions équivalentes :

\begin{enumerate}
\item[$\alpha$)] $X \to \widetilde{X}$ est surjective ;
\item[$\beta$)] toute partie fermée filtrante de $X$ a un plus grand élément ;
\item[$\gamma$)] toute suite croissante de $X$ est stationnaire ;
\item[$\delta$)] tout point de $\widetilde{X}$ est essentiel.
\end{enumerate}

$\alpha$) et $\beta$) sont deux façons de dire la même chose. L'équivalence
avec $\gamma$) tient à ce qu'un ensemble filtrant sans élément maximal
contient une suite strictement croissante, et qu'un ensemble filtrant avec un
élément maximal a celui-ci pour plus grand élément.\footnote{Ce passage de la
condition de chaîne \emph{séquentielle} à la condition sur les parties
filtrantes quelconques repose sur l'axiome du choix dépendant ; la page ne le
mentionne pas. Elle écrit d'ailleurs « artinien » et « suite décroissante »,
qui est la même condition lue dans sa convention d'ordre : le sens de
l'inégalité une fois fixé comme ici, c'est la condition de chaîne
\emph{ascendante}.} Quant à $\delta$), il repose sur ce que la page note en
marge, d'un trait et d'un point d'exclamation : \emph{l'image de $X$ dans
$\widetilde{X}$ est exactement l'ensemble des points essentiels de
$\widetilde{X}$}.

Un \emph{point essentiel} d'un topos est un point dont le foncteur image
inverse — le foncteur fibre — admet lui-même un adjoint à gauche ; pour un
espace, c'est un point qui a un plus petit voisinage ouvert, et c'est encore un
point dont le foncteur fibre commute aux produits quelconques.\footnote{Cette
troisième formulation est celle que donne la page 48. Nous
la rapportons ici parce qu'elle est la définition la plus maniable et qu'elle
éclaire le mot « essentiel » : un morphisme de topos est dit essentiel lorsque
son image inverse a un adjoint à gauche, exactement comme à la page 15.}

La page 26 clôt le dictionnaire par sa forme la plus forte. Pour un topos $E$,
sont équivalentes :

\begin{enumerate}
\item[$\alpha$)] $E$ est équivalent à la fois à un $\operatorname{Top}(X)$ et à
un $\widehat{C}$ ;
\item[$\beta$)] $E \simeq \operatorname{Top}(X)$ avec $X$ d'Alexandrov ;
\item[$\gamma$)] $E \simeq \widehat{I}$ avec $I$ un ensemble ordonné ;
\item[$\delta$)] $E$ est engendré par les sous-objets de son objet final, et a
suffisamment de points essentiels.
\end{enumerate}

Un tel $E$ est appelé \emph{espace topologique essentiel}, et l'on a des
équivalences de 2-catégories entre les topos spatiaux essentiels, les
préordres et les espaces d'Alexandrov. Deux tests accompagnent l'énoncé :
$\widehat{C}$ est de ce type si et seulement si $C$ est préordonnée ;
$\operatorname{Top}(X)$ l'est si et seulement si l'ensemble des points
essentiels de $X$ est \emph{très dense}.\footnote{Très dense au sens de
SGA 4 : l'inclusion induit une bijection entre les ouverts de $X$ et ceux du
sous-espace, donc une équivalence des topos. C'est exactement la condition
qu'il faut, et c'est elle que la page 48 reprendra pour répondre à la
question du bas de la page 40.}

\pagerange{28}{28}
\section*{Le produit, du côté des treillis (page 28)}

Un feuillet isolé, dont la moitié supérieure porte un calcul de quaternions et
de formules de duplication trigonométriques sans aucun rapport — le dossier est
fait de papiers réemployés. Le bas de la page revient au produit des pages 9 à
12, et en donne la face manquante.

Soient $\sigma$ un treillis à bornes supérieures quelconques et bornes
inférieures finies distributives par rapport à elles — un \emph{cadre}, dirait-on
aujourd'hui — et $\alpha : \mathcal{O}_{X} \to \sigma$,
$\beta : \mathcal{O}_{Y} \to \sigma$ deux morphismes de cadres. La question est
de savoir si le couple $(\alpha,\beta)$ provient d'un unique morphisme
$\gamma$ défini sur les ouverts de $X \times Y$.

L'injectivité est immédiate : $\gamma$ est connu dès que l'on connaît
$\gamma(U \times V) = \alpha(U) \wedge \beta(V)$, puisque tout ouvert de
$X \times Y$ est borne supérieure de rectangles. Pour la surjectivité, on pose
\[
\gamma(W) \;=\; \operatorname*{Sup}_{U \times V \subset W}
\alpha(U) \wedge \beta(V),
\]
et tout revient à vérifier que $\gamma$ commute aux bornes inférieures
finies. Le calcul de la page est complet et correct :
\[
\gamma(W) \wedge \gamma(W')
= \operatorname*{Sup}_{\substack{U \times V \subset W \\ U' \times V' \subset W'}}
\alpha(U \cap U') \wedge \beta(V \cap V')
= \gamma(W \cap W'),
\]
la première égalité par distributivité dans $\sigma$ et parce que $\alpha$,
$\beta$ commutent aux intersections finies, la seconde parce que
$(U \cap U') \times (V \cap V')$ est bien contenu dans $W \cap W'$ et que tout
rectangle contenu dans $W \cap W'$ est de cette forme. Reste
$\gamma(X \times Y) = \alpha(X) \wedge \beta(Y) = 1$, et le calcul
$\gamma(X' \times Y) = \alpha(X')$ que la page mène en supposant
$Y \neq \emptyset$.

C'est donc oui, et l'objet ainsi caractérisé est le \emph{produit de
locales} $\mathcal{O}_{X} \otimes \mathcal{O}_{Y}$ : le treillis des ouverts du
produit \emph{de topos}, celui que les pages 9 à 12 appelaient la topologie
$\pi$.\footnote{La page ne fait pas ce rapprochement, et les deux feuillets
sont éloignés dans le dossier. Il est pourtant exact, et il fournit la
justification que la démonstration des pages 9 et 10 utilisait sans la
démontrer : $W \mapsto W \times V$ commute aux bornes supérieures quelconques,
ce qui est ici la définition même de $\gamma$.}


\pagerange{30}{30}
\section*{Un lemme sur les algèbres d'Azumaya (pages 30 à 32)}

Ce lemme n'a rien à voir avec le reste du dossier, et il faut le prendre pour
ce qu'il est : une page d'algèbre, rangée là par accident de classement.

Soit $E$ un topos localement annelé. On considère deux catégories fibrées sur
$E$ :

\begin{itemize}
\item $C$, dont les objets sont les couples $(\mathcal{G}, \mathcal{E})$ de
modules localement libres de type fini, partout non nuls ;
\item $C'$, dont les objets sont les quadruplets
$(A, \mathcal{E}, \mathcal{F}, \varphi)$ où $A$ est une algèbre localement
isomorphe à une algèbre de matrices — une \emph{algèbre d'Azumaya} —, où
$\mathcal{E}$ et $\mathcal{F}$ sont localement libres de type fini et partout
non nuls, et où $\varphi$ est un isomorphisme
\[
A \otimes \underline{\operatorname{End}}(\mathcal{E}) \;\simeq\;
\underline{\operatorname{End}}(\mathcal{F}).
\]
\end{itemize}

\textbf{Lemme.} \emph{Le foncteur}
\[
F : (\mathcal{G}, \mathcal{E}) \longmapsto
\bigl(\underline{\operatorname{End}}(\mathcal{G}),\ \mathcal{E},\
\mathcal{G} \otimes \mathcal{E},\ \varphi_{\mathcal{G},\mathcal{E}}\bigr)
\]
\emph{— où $\varphi_{\mathcal{G},\mathcal{E}}$ est l'isomorphisme canonique
$\underline{\operatorname{End}}(\mathcal{G}) \otimes
\underline{\operatorname{End}}(\mathcal{E}) \simeq
\underline{\operatorname{End}}(\mathcal{G} \otimes \mathcal{E})$ — est une
équivalence.}\footnote{Ce que dit ce lemme, en langage d'aujourd'hui : la
donnée de $\varphi$ \emph{est} une trivialisation de la classe de $A$ dans le
groupe de Brauer de $E$, et le lemme affirme qu'une telle trivialisation
équivaut exactement à la donnée du module $\mathcal{G}$ qui scinde $A$ — ni
plus, ni moins. La théorie du groupe de Brauer d'un schéma est celle des trois
exposés \emph{Le groupe de Brauer} (1966--1968) ; les gerbes qui la portent
sont celles de Giraud (1971). Les pages ne sont pas datées.}

\pagerange{30}{32}
\subsection*{La démonstration}

Elle procède par dévissage sur les rangs, ce qui est l'idée juste : $F$
provient d'un morphisme de champs $\underline{C} \to \underline{C}'$, et il
suffit de prouver que celui-ci est une équivalence. En décomposant selon les
rangs,
\[
\underline{C} \simeq \coprod_{g,e} \underline{C}_{g,e},
\qquad
\underline{C}' \simeq \coprod_{a,e,f} \underline{C}'_{a,e,f},
\]
le foncteur induit est $F_{g,e} : \underline{C}_{g,e} \to
\underline{C}'_{g^{2},\,e,\,ge}$.\footnote{Le troisième rang est illisible sur
la page ; il vaut nécessairement $ge$, puisque $\mathcal{F} = \mathcal{G}
\otimes \mathcal{E}$. La convention de rang pour $A$ est celle de la page :
$A$ de rang $a^{2}$, donc $a = g$ ici.} On est ramené à prouver que chaque
$F_{g,e}$ est une équivalence, et comme il s'agit de champs en groupoïdes,
cela se scinde en deux vérifications : deux objets sont localement isomorphes,
et $F$ est bijectif sur les automorphismes.

La première est le point où le vocabulaire hésite. La page écrit
« $\underline{C}_{g,e}$ est une gerbe (liée par $Gl(e) \times Gl(g)$) », puis
barre entièrement les deux lignes et reprend : « c'est essentiellement un champ
en groupoïdes, donc reste à prouver que deux objets sont localement
isomorphes ».\footnote{C'est le mot « gerbe » qui est abandonné, non l'idée. Un
champ en groupoïdes localement non vide dont deux objets sont localement
isomorphes \emph{est} une gerbe, et c'est bien de cela qu'il s'agit ; le lien
est celui que la ligne barrée nomme. La rature date sans doute d'une époque où
le terme n'était pas encore stabilisé, ce qui est un des rares indices
chronologiques du dossier.} La vérification elle-même est courte : deux
isomorphismes $\varphi, \varphi'$ diffèrent par un automorphisme de
$\underline{\operatorname{End}}(\mathcal{F})$, lequel est localement intérieur,
donc provient localement d'un automorphisme de $\mathcal{F}$.

Pour la seconde, un automorphisme $(u,v)$ de $\xi = (\mathcal{G}, \mathcal{E})$
a pour image $\bigl(\underline{\operatorname{End}}(u),\ v,\ u \otimes v\bigr)$.
L'injectivité est immédiate : $v$ est lu directement, $u \otimes v$ aussi,
d'où $u$ puisque $\mathcal{E}$ est partout non nul. Pour la surjectivité, on
part d'un automorphisme $\tau = (\lambda, v, w)$ de $\xi'$, dont la
compatibilité s'écrit

\begin{tikzcd}[column sep=large]
\underline{\operatorname{End}}(\mathcal{G}) \otimes \underline{\operatorname{End}}(\mathcal{E}) \arrow[r, "\varphi"] \arrow[d, "\lambda \otimes \underline{\operatorname{End}}(v)"'] & \underline{\operatorname{End}}(\mathcal{G} \otimes \mathcal{E}) \arrow[d, "\underline{\operatorname{End}}(w)"] \\
\underline{\operatorname{End}}(\mathcal{G}) \otimes \underline{\operatorname{End}}(\mathcal{E}) \arrow[r, "\varphi"'] & \underline{\operatorname{End}}(\mathcal{G} \otimes \mathcal{E})
\end{tikzcd}

et il s'agit de trouver localement $u$ avec $\lambda =
\underline{\operatorname{End}}(u)$ et $w = u \otimes v$. Le premier point est
encore le caractère localement intérieur de $\lambda$ ; quitte à composer par
$F(u,v)^{-1}$, on peut donc supposer $\lambda = \mathrm{id}$, et il reste à
voir que $w = \mathrm{id}_{\mathcal{G}} \otimes v$. La commutativité donne
\[
w\,(\mathrm{id}_{\mathcal{G}} \otimes X)\,w^{-1}
= \mathrm{id}_{\mathcal{G}} \otimes v X v^{-1}
\qquad \text{pour tout endomorphisme } X \text{ de } \mathcal{E},
\]
c'est-à-dire que $w\,(\mathrm{id} \otimes v)^{-1}$ centralise le facteur
$\mathrm{id}_{\mathcal{G}} \otimes \underline{\operatorname{End}}(\mathcal{E})$
— et la page devient ici trop surchargée pour être suivie.\footnote{Le pas
manquant est le lemme du double centralisateur : le centralisateur de
$1 \otimes \underline{\operatorname{End}}(\mathcal{E})$ dans
$\underline{\operatorname{End}}(\mathcal{G} \otimes \mathcal{E})$ est
$\underline{\operatorname{End}}(\mathcal{G}) \otimes 1$, d'où
$w = u \otimes v$ pour un $u$ localement défini. Nous le signalons sans
l'attribuer à la page, qui s'interrompt avant. La page 31, au crayon et restée
à l'état d'esquisse, porte des torseurs $\mathrm{TORS}^{n}_{U}(X,\Pi)$ et
l'isomorphisme $\pi_{0}(T^{n}) \simeq H^{n}(X,\Pi)$ : c'est la cohomologie non
abélienne dans laquelle ce lemme s'inscrit, mais rien n'y est assez formé pour
être restitué.}


\pagerange{35}{35}
\section*{Espèces de structure et types de limites (pages 35 à 39)}

L'autre manuscrit, celui des pages impaires. Il prend les \emph{théories}
elles-mêmes pour objets, et sa question finale est de savoir si les corps en
forment une.

\pagerange{35}{35}
\subsection*{Une espèce de structure comme foncteur}

Soit $T$ une \emph{espèce de structure} définie par limites projectives finies
(resp. quelconques) sur une famille d'ensembles de base $(E_{i})_{i \in I}$,
$I$ fixé. Pour toute catégorie $C$ à limites finies, on dispose de la catégorie
$T(C)$ des structures d'espèce $T$ dans $C$, et du foncteur « famille d'objets
sous-jacente »
\[
\varphi^{T}_{C} : T(C) \longrightarrow C^{I} ;
\]
pour tout foncteur $u : C \to C'$ commutant aux limites finies, un
$T(u) : T(C) \to T(C')$ au-dessus de $u^{I}$. Ces données étant fonctorielles,
$T$ est un foncteur
\[
T : L_{0} = (\mathrm{Cat} \text{ à } \varprojlim \text{ finies})
\longrightarrow \mathrm{Cat},
\]
d'où une catégorie fibrée $\mathcal{T}$ munie de $\varphi^{T} : \mathcal{T} \to
\mathcal{T}_{I}$.\footnote{C'est la présentation des \emph{théories
essentiellement algébriques}, ou théories à limites : une structure est une
famille d'objets munie d'opérations dont les domaines sont définis par des
équations, et une telle espèce se laisse évaluer dans n'importe quelle
catégorie à limites finies. Les esquisses d'Ehresmann et les théories de
Lawvere en sont les deux formulations contemporaines.} Pour deux espèces $T$,
$T'$ on définit $\underline{\operatorname{Hom}}_{L}(T,T')$, et si elles portent
le même $I$, la catégorie $\underline{\operatorname{Hom}}_{L}(I;T,T')$ des
homomorphismes munis d'une donnée de compatibilité entre $\varphi^{T}$ et
$\varphi^{T'}$.

\pagerange{37}{37}
\subsection*{Quand la théorie est représentable}

La page 37 pose la condition qui donne à tout cela sa forme définitive.
\emph{Supposons $T$ 2-représentable} : il existe $\beta^{T}$, catégorie à
limites, telle que
\[
T(C) \;\simeq\; \underline{\operatorname{Hom}}_{\varprojlim}(\beta^{T},\, C)
\qquad \text{pour toute } C \text{ à limites.}
\]
Alors trois choses suivent.

\begin{enumerate}
\item[a)] $\beta^{T}$ est unique à équivalence près, et l'équivalence unique à
isomorphisme unique près.
\item[b)] Pour toute catégorie 2-cofibrée $\mathcal{F}$ sur
$(\mathrm{Cat}\ \varprojlim)$, on a
$\underline{\operatorname{Hom}}(T, \mathcal{F}) \simeq \mathcal{F}(\beta)$ ; en
particulier, si $\mathcal{F}$ est définie par une espèce $T'$,
\[
\underline{\operatorname{Hom}}(T, T') \;\simeq\; T'(\beta).
\]
\item[c)] En prenant pour $T'$ l'espèce « ensemble de base sans structure »,
pour laquelle $T'(\beta) = \beta$, on trouve
\[
\beta \;=\; \underline{\operatorname{Hom}}(T,\, i) \;=\; \Gamma^{T}.
\]
\end{enumerate}

Ce dernier énoncé est le cœur du fragment, et il vaut d'être dit en clair :
\emph{la théorie se retrouve comme la catégorie de ses propres opérations
naturelles}. $\Gamma^{T}$ est ce que toute structure d'espèce $T$ sait
fabriquer, uniformément en la structure ; l'énoncé dit que cela suffit à
reconstituer la théorie.\footnote{C'est, sous une forme antérieure au
vocabulaire, la dualité qu'on appelle aujourd'hui \emph{de Gabriel--Ulmer}
(1971) : entre les petites catégories à limites finies et les catégories
localement présentables, une théorie et sa catégorie de modèles se déterminent
l'une l'autre. La même construction, avec la même notation $\Gamma^{T}$ et le
même « ensemble de base sans structure » pour objet dualisant, occupe la
page 13 de la cote 161-1, où elle est écrite indépendamment. L'énoncé
sera poussé jusqu'à sa forme explicite aux pages 45 et 47.}

\pagerange{39}{39}
\subsection*{Les corps forment-ils une espèce de structure ? (page 39)}

Vient enfin la question qui occupe la fin du manuscrit, et qui est posée avec
une netteté remarquable.

L'espèce de structure « anneau commutatif » est définie par limites finies ;
notons $\lambda_{0}$ le type correspondant. Fixons pour tout le dossier la
notation : un \emph{type} $\lambda$ est un couple
$\lambda = (\underleftarrow{\lambda}, \underrightarrow{\lambda})$ formé d'une
classe de types de diagrammes dont on exige les limites et d'une classe dont on
exige les colimites ; $\mathrm{Cat}_{\lambda}$ est la 2-catégorie des
catégories qui les possèdent et des foncteurs qui y commutent ; et tout
$\lambda$ recevant $\lambda_{0}$ définit une espèce $T_{\lambda}$.

\textbf{Question.} \emph{Pour quels $\lambda$ existe-t-il un sous-type plein
$T^{*}_{\lambda}$ — c'est-à-dire $T^{*}_{\lambda}(C) \to T_{\lambda}(C)$
pleinement fidèle pour tout $C$ — tel que, pour $C = (\mathrm{Ens})$, on
retrouve exactement la catégorie des corps commutatifs ? Et un tel sous-type
est-il unique ?}

La machinerie de réponse est en place dès cette page. $T_{\lambda}$ est
$\lambda$-représentable par la sous-catégorie pleine $S^{T}_{\lambda}$ de
$(\mathrm{Ann})$ engendrée, par les colimites de type $\lambda$, par l'objet
libre à un générateur $\mathbb{Z}[t]$ ; et les sous-types pleins de
$T_{\lambda}$ correspondent aux \emph{catégories de fractions} de
$S^{T}_{\lambda}$ par un ensemble $M$ de flèches stable par composition et par
les colimites de type $\lambda$. La traduction de la condition « les modèles
ensemblistes sont les corps » est alors :
\[
M \;=\; \bigl\{\, u : A \to B \ \bigm|\ \operatorname{Hom}(B,k)
\xrightarrow{\ \sim\ } \operatorname{Hom}(A,k) \ \text{pour tout corps } k
\,\bigr\},
\]
c'est-à-dire, géométriquement, les morphismes de schémas affines
$X \to Y$ tels que $X(k) \simeq Y(k)$ pour tout corps $k$ — les morphismes
bijectifs à extensions résiduelles triviales. Cet ensemble de flèches est
manifestement stable par composition et par colimites de type $\lambda$, et il
le serait pour n'importe quelle sous-catégorie de $(\mathrm{Ann})$.

L'unicité est acquise dès que $\underrightarrow{\lambda} = \emptyset$. Reste
l'existence, et la page ouvre l'énumération des cas : \emph{a) si
$\lambda = \lambda_{0}$, il n'existe pas de tel sous-type} — la démonstration
est page 41, dont la première ligne achève la phrase que
celle-ci laisse en suspens.\footnote{C'est une continuité facile à
manquer : la page 40 s'intercale
physiquement entre les pages 39 et 41 sans appartenir au même manuscrit. La
lecture qui suit les feuillets dans l'ordre y perd le fil de la phrase ; celle
qui suit les fils le garde.}

\pagerange{41}{41}
\section*{Les corps sont-ils une espèce de structure ? (page 41)}

Cette page achève la phrase commencée au bas de la page 39, et il faut la lire dans son prolongement immédiat : la page 40 qui
les sépare physiquement appartient à l'autre manuscrit.\footnote{C'est aussi la
page la plus difficile du dossier : très cursive, lourdement biffée, elle porte
à elle seule le tiers des passages illisibles du dossier. Nous restituons ci-dessous
la charpente de l'argument \emph{autour} des lacunes, sans rien reconstruire de
ce qu'elles cachent ; là où un pas manque, la note le dit. En particulier
l'hypothèse de tête, notée $\lambda = \mathcal{E}\lambda_{0}$ sur la page, est
donnée par la transcription comme une lecture incertaine, et nous ne bâtissons
rien sur elle.}

Rappelons la mise en place. Un \emph{type} est un couple
$\lambda = (\underleftarrow{\lambda}, \underrightarrow{\lambda})$ — les
diagrammes dont on exige les limites, ceux dont on exige les colimites ; $T$
est l'espèce de structure « anneau commutatif », de type $\lambda_{0}$ ;
$S^{T}_{\lambda}$ est la sous-catégorie pleine de $(\mathrm{Ann})$ engendrée par
$\mathbb{Z}[t]$ sous les colimites de type $\lambda$ ; et l'on cherche un
sous-type plein $T^{*}_{\lambda}$ dont les modèles ensemblistes soient les
corps. Le critère est celui qu'énonce la première phrase de la page :

\begin{quote}
La question a une solution affirmative si et seulement si tout anneau
$\mathcal{C}$ qui rend bijective l'application
$\operatorname{Hom}(B,\mathcal{C}) \to \operatorname{Hom}(A,\mathcal{C})$ pour
toute flèche $u : A \to B$ de $M$, est un corps.
\end{quote}

$M$ désignant, comme à la page 39, l'ensemble des flèches que tout corps
inverse.

\pagerange{41}{41}
\subsection*{a) Avec les seuls moyens qui définissent les anneaux : non}

Si $\lambda = \lambda_{0}$, la catégorie $S^{T}_{\lambda}$ se réduit aux
anneaux de polynômes $\mathbb{Z}[t_{1}, \ldots, t_{n}]$, et une flèche entre
deux d'entre eux qui induit une bijection sur les $k$-points pour \emph{tout}
corps $k$ est un isomorphisme — la page invoque là le théorème principal de
Zariski.\footnote{L'identification de $S^{T}_{\lambda_{0}}$ aux anneaux de
polynômes est la nôtre ; la page ne l'écrit pas. L'argument intermédiaire — une
demi-douzaine de mots dont « fibrée par fibre » et « bien fermée », plusieurs
illisibles — n'est pas restituable, et nous ne le restituons pas. Seuls sont
lisibles l'appel à Zariski et la conclusion, qui suffisent à porter l'énoncé.}
Donc $M$ est réduit aux isomorphismes ; tout anneau inverse $M$ ; le critère
demanderait que tout anneau soit un corps. Absurde : il n'existe pas de tel
sous-type.

\pagerange{41}{41}
\subsection*{b) Avec la localisation en plus : oui}

Supposons maintenant que $S^{T}_{\lambda}$ contienne, à côté de
$\mathbb{Z}[t]$, l'anneau $\mathbb{Z}[t,t^{-1}]$ — ce qui a lieu dès que
$\lambda$ autorise toutes les limites finies de schémas affines, auquel cas
$S^{T}_{\lambda}$ est la catégorie des anneaux de type fini sur
$\mathbb{Z}$.\footnote{La page écrit cette hypothèse sous la forme
$\lambda = (\underleftarrow{\Sigma}, \underrightarrow{\lambda} = \emptyset)$,
lecture que la transcription donne pour douteuse. L'hypothèse substantielle,
elle, est parfaitement lisible et c'est la seule dont l'argument se sert :
\emph{$\mathbb{Z}[t,t^{-1}]$ est disponible}. Nous énonçons celle-là.
$\underrightarrow{\lambda} = \emptyset$ assure par ailleurs l'unicité du
sous-type, acquise dès la page 39.} Considérons alors la flèche canonique
\[
u_{0} : \mathbb{Z}[t] \longrightarrow \mathbb{Z}[t,t^{-1}] \times \mathbb{Z},
\qquad t \longmapsto (t,\,0),
\]
et notons que le but est de présentation finie :
$\mathbb{Z}[t,t^{-1}] \times \mathbb{Z} \simeq \mathbb{Z}[x,y] \big/
\bigl(x(xy-1),\, y(xy-1)\bigr)$, l'idempotent $e = xy$ découpant les deux
facteurs.\footnote{Vérification : dans le quotient, $e = xy$ vérifie
$e^{2} = e$ ; sur le facteur $e = 1$ on a $xy = 1$, d'où
$\mathbb{Z}[x,x^{-1}]$ ; sur le facteur $e = 0$ les deux relations donnent
$x = y = 0$, d'où $\mathbb{Z}$. La page écrit d'abord ce quotient avec un autre
système de générateurs, qu'elle raye.}

Appliquons $\operatorname{Hom}(-,\mathcal{C})$. Un homomorphisme d'un produit
d'anneaux dans $\mathcal{C}$ équivaut à une décomposition
$\mathcal{C} = \mathcal{C}_{1} \times \mathcal{C}_{2}$ — c'est-à-dire à un
idempotent, c'est-à-dire à une partition de $\operatorname{Spec}\mathcal{C}$ en
deux ouverts-fermés — accompagnée d'un élément inversible
$\varepsilon \in \mathcal{C}_{1}^{\times}$ ; et $u_{0}$ l'envoie sur l'élément
$(\varepsilon, 0)$ de $\mathcal{C}$. La transformée de $u_{0}$ est donc
\[
\coprod_{\mathcal{C} = \mathcal{C}_{1} \times \mathcal{C}_{2}}
\mathcal{C}_{1}^{\times} \longrightarrow \mathcal{C},
\]
et \emph{elle est bijective si et seulement si $\mathcal{C}$ est un corps} :
tout élément doit s'écrire, d'une seule manière, comme inversible sur une pièce
et nul sur la pièce complémentaire — ce qui force $\operatorname{Spec}
\mathcal{C}$ connexe, c'est-à-dire $\mathcal{C}$ sans idempotent non trivial, et
tout élément inversible ou nul.\footnote{Il faut exclure à part l'anneau nul,
pour lequel l'application est encore bijective ; la convention moderne veut
qu'un corps soit non nul, et la page raisonne dans cette convention sans la
dire. La page s'interrompt d'ailleurs au milieu de la justification, le crochet
ouvert restant ouvert : elle a écrit la moitié « connexe » et non la moitié
« tout élément inversible ou nul ». Les deux sont nécessaires, et les deux
suivent de la bijectivité.}

En particulier $u_{0}$ appartient à $M$, et à lui seul il force le critère :
tout anneau inversant $M$ est un corps. \emph{Le sous-type existe.} La réponse
du dossier n'est donc pas un refus mais une ligne de partage — les corps
échappent aux moyens qui définissent les anneaux, et sont récupérés dès qu'on
ajoute la localisation.\footnote{C'est, mot pour mot, la distinction moderne
entre une théorie \emph{essentiellement algébrique} — définie par limites
finies, dont les modèles sont stables par produits, ce que la classe des corps
n'est pas — et une théorie \emph{géométrique}, où l'axiome des corps s'écrit
précisément sous la forme trouvée ici : le morphisme
$\mathcal{C}^{\times} + \{0\} \to \mathcal{C}$ est un isomorphisme. C'est
l'axiome sous lequel les corps apparaissent dans le topos de Zariski, chez
Hakim (\emph{Topos annelés et schémas relatifs}, 1972). Les pages ne sont pas
datées, et rien n'indique de quel côté de cette date elles tombent.}

\pagerange{43}{43}
\section*{Catégories $\lambda$-abstraites (pages 43, 45, 47)}

Les trois pages impaires suivantes reprennent la machinerie des pages 35 et 37 —
une espèce de structure $T$, la catégorie $T(C)$ de ses modèles dans $C$, le
foncteur d'oubli $b^{T}_{C} : T(C) \to C^{I}$ — et la poussent jusqu'à un
énoncé de reconstitution.

\pagerange{43}{43}
\subsection*{Une page de travail}

La page 43 est un feuillet de notes : le carré comparant $\mathrm{Cat}_{\lambda}$
et $\mathrm{Cat}_{\lambda'}$ pour $\lambda \subset \lambda'$, le carré
2-cartésien
\[
\operatorname{Hom}_{\lambda}(R, C) \hookrightarrow
\operatorname{Hom}_{\lambda}(R, \widehat{C}\,)
\quad\text{au-dessus de}\quad
C^{I} \hookrightarrow \widehat{C}^{\,I},
\]
et une liste de cinq questions ouvertes qu'il vaut de garder telle quelle :
rigidité de $R_{T}$ ; existence des colimites de type $\lambda$ dans $T(C)$ ;
existence des limites ; existence d'un adjoint au foncteur d'oubli
$\varphi : T(\mathrm{Ens}) \to (\mathrm{Ens})^{I}$ lorsque $\lambda$ est formé
de diagrammes finis ; et enfin
\[
R_{T} \;\overset{?}{\simeq}\; \operatorname{Hom}_{\lambda}
\bigl((T(\mathrm{Ens}))_{\circ},\, (\mathrm{Ens})\bigr),
\]
avec deux points d'interrogation de sa main.\footnote{La quatrième question a
aujourd'hui une réponse : pour une théorie à limites finies, le foncteur
d'oubli d'une catégorie localement présentable admet un adjoint à gauche —
l'objet libre existe. La cinquième est la forme précise de la dualité qui suit.}

\pagerange{45}{45}
\subsection*{La reconstitution}

La page 45 pose les hypothèses. Soit $T$ une catégorie $\lambda$-abstraite sur
$\mathrm{Cat}_{\lambda}$, munie de ses foncteurs $b^{T}_{C} : T(C) \to C^{I}$,
et supposons :

\begin{enumerate}
\item[a)] pour tout foncteur pleinement fidèle $C \to C'$ dans
$\mathrm{Cat}_{\lambda}$, $T(i) : T(C) \to T(C')$ est pleinement fidèle, et le
carré qu'il forme avec les $b^{T}$ est 2-cartésien — il suffit de l'exiger pour
$C \hookrightarrow \widehat{C}$ ;
\item[b)] pour tout $C$, le foncteur canonique
\[
T(\widehat{C}\,) \longrightarrow \operatorname{Hom}_{\lambda}(C^{\circ}, S),
\qquad S \overset{\text{déf}}{=} T(\mathrm{Ens}),
\]
est une équivalence de catégories.\footnote{La page note en marge, et c'est une
mise en garde qu'il faut garder : « $\widehat{C}$ n'est pas une
$\mathcal{U}$-catégorie ! » Les énoncés qui suivent supposent qu'on a choisi un
univers assez grand, ou qu'on s'est restreint aux petites catégories ; la
question de taille est réelle et le manuscrit ne la tranche pas.}
\end{enumerate}

Alors $T(C)$ se reconstitue à partir de deux données seulement : la catégorie
$S$ des modèles ensemblistes, et la famille $\beta = (\beta_{i})_{i \in I}$ des
foncteurs fibres $\beta_{i} : S \to (\mathrm{Ens})$. Précisément,
\[
T(C) \;=\; \bigl\{\, \varphi : C^{\circ} \to S \ \text{dans}
\ \operatorname{Hom}_{\lambda} \ \bigm| \ \beta_{i} \circ \varphi
\ \text{représentable pour tout } i \,\bigr\},
\]
sous-catégorie pleine de $\operatorname{Hom}_{\lambda}(C^{\circ},
S)$.\footnote{Une question reste ouverte en marge, et elle est la bonne : « si
$f : C \to C'$ est un $\lambda$-foncteur, comment définir $T(f)$ en fonction de
$S$ et des $(\beta_{i})$ ? » La reconstitution est donnée objet par objet, non
comme reconstitution du foncteur $T$ tout entier.}

Reste à identifier la théorie elle-même. Soit $R$ la sous-catégorie pleine de
$\operatorname{Hom}_{\lambda}(S, (\mathrm{Ens}))$ formée des $\psi$ tels que
$\psi \circ \xi$ soit représentable pour tout $C$ et tout $\xi \in T(C)$ : elle
est stable par les limites de type $\underleftarrow{\lambda}$, et l'on dispose
d'un accouplement $T(C) \times R \to C$, d'où
\[
T(C) \longrightarrow \operatorname{Hom}_{\lambda}(R, C).
\]
\emph{À quelle condition est-ce une équivalence ?} Deux conditions sont
nécessaires, et ce sont elles qui donnent sa forme finale à l'énoncé :

\begin{enumerate}
\item[a)] $R$ doit être formé de foncteurs représentables, c'est-à-dire
$R \simeq \Sigma^{\circ}$ pour une sous-catégorie pleine $\Sigma \subset S$ —
la donnée de $R$ équivaut alors à celle de $\Sigma$ ;
\item[b)] $\Sigma$ doit engendrer $S$, au sens où
$S \to \operatorname{Hom}(\Sigma^{\circ}, (\mathrm{Ens})) = \widehat{\Sigma}$
est pleinement fidèle.
\end{enumerate}

\pagerange{47}{47}
\subsection*{La forme finale}

La page 47 prend $\lambda$ = l'ensemble des types de diagrammes \emph{finis} et
énonce le résultat dans l'autre sens : on part de $\Sigma$ et l'on fabrique
$S$. Soient donc $\Sigma$ une catégorie stable par limites de type $\lambda$,
$R = \Sigma^{\circ}$, et
\[
S \;=\; \operatorname{Hom}_{\underrightarrow{\lambda}}(R, \mathrm{Ens}),
\qquad \Sigma \;\subset\; S \;\subset\; \widehat{\Sigma},
\]
les foncteurs sur $R$ commutant aux colimites de type
$\underrightarrow{\lambda}$.\footnote{La page écrit la chaîne d'inclusions en
sens inverse, $S \supset \widehat{\Sigma} \supset \Sigma$. C'est un lapsus :
les foncteurs astreints à commuter à des colimites forment une sous-catégorie
pleine de $\widehat{\Sigma}$, laquelle contient les représentables. La ligne
porte d'ailleurs un « $\mathrm{Ens}$ » rayé, signe d'une écriture hésitante.}
Alors :

\begin{enumerate}
\item[a)] $S$ admet des colimites quelconques ;
\item[b)] pour toute catégorie $\mathcal{T}$ à colimites quelconques, la
restriction
\[
\operatorname{Hom}_{\underrightarrow{\lambda}}(\Sigma, \mathcal{T})
\xrightarrow{\ \sim\ } \operatorname{Hom}_{\lambda'}(S, \mathcal{T})
\]
est une équivalence, $\operatorname{Hom}_{\lambda'}$ désignant les foncteurs
cocontinus ;
\item[c)] $S$ est engendré par $\Sigma$ sous les colimites.
\end{enumerate}

Autrement dit : $S$ est la cocomplétion libre de $\Sigma$ \emph{respectant les
colimites que $\Sigma$ possède déjà}, et $\Sigma$ s'y retrouve comme la
sous-catégorie des objets de présentation finie. La théorie $R = \Sigma^{\circ}$
et la catégorie de modèles $S$ se déterminent donc l'une l'autre, et la
correspondance est celle des propriétés universelles.\footnote{C'est la
\emph{dualité de Gabriel--Ulmer} (1971), entre les petites catégories à limites
finies et les catégories localement de présentation finie, dont ces deux pages
énoncent les deux moitiés : la page 45 va des modèles à la théorie, la page 47
de la théorie aux modèles. Le $\lambda$ du manuscrit est un ensemble de
\emph{types de diagrammes} ; le $\lambda$ d'aujourd'hui est un cardinal
régulier, et l'on parle de catégories localement $\lambda$-présentables — même
lettre, même rôle, mais l'index a changé de nature. La page 37 avait
donné la moitié abstraite du même énoncé, $\beta = \Gamma^{T}$.}


\pagerange{34}{34}
\section*{Notes de cours : les topos contre les espaces (pages 34 à 40)}

À partir d'ici, deux manuscrits alternent. Les pages paires — 34, 36, 38, 40,
puis 42, 44, 46 et 48 — portent des notes de cours sur les
topos ; les pages impaires, un tout autre sujet, repris plus bas.

Ces pages-ci sont écrites \emph{en anglais}. Ce qui suit n'en est pas une
traduction mais une restatement en français, comme tout le reste de cette
édition ; qui veut les mots de la page ira à la transcription, qui suit la
langue du feuillet.\footnote{C'est le seul endroit du dossier où cette édition
s'écarte de la page par la langue, et il fallait le décider une fois pour
toutes : les pages 34, 36, 38, 40, 48, 52, 53 et 54 sont en anglais. Nous ne conservons de l'anglais que les tournures
citées entre guillemets, quand c'est la formule elle-même qui compte — « algèbre
$=$ topologie ! » en est une.}

\pagerange{34}{34}
\subsection*{Ce que $\operatorname{Top}(X)$ retient de $X$}

Le cours part d'un rappel — images directes et inverses, foncteur fibre comme
image inverse, faisceaux à structures algébriques, définis par limites finies
et préservés en sens contraire par les deux images — puis pose sa thèse :
\emph{la catégorie $\operatorname{Top}(X)$ contient toute l'information
topologique exploitable sur $X$}. La liste qui l'appuie vaut d'être gardée
telle quelle : les faisceaux à structures algébriques, la notion de champ sur
$X$, la connexité, la cohomologie et l'homotopie, le rapport entre $\pi_{1}$ et
les revêtements, l'algèbre homologique non commutative de $H^{1}$, la
connexité locale, la totale discontinuité, la quasi-compacité, et les
opérations sur les espaces (sommes, sommes amalgamées) telles que les
opérations sur les topos les reflètent.

Le point technique est le suivant. Notons $\mathrm{Fib}(E)$ l'ensemble des
classes d'isomorphie de foncteurs fibres généralisés. Pour
$E = \operatorname{Top}(X)$, il s'identifie à l'ensemble $\widetilde{X}$ des
fermés irréductibles de $X$, muni de la topologie évidente, et l'application
canonique est $x \mapsto \overline{\{x\}}$. D'où trois équivalences, que la page
énonce avec la précision qui manquait à la page 24 :

\begin{itemize}
\item elle est injective si et seulement si $\overline{\{x\}} =
\overline{\{y\}} \Rightarrow x = y$ — l'axiome $T_{0}$ ;
\item elle est surjective si et seulement si tout fermé irréductible a un point
générique ;
\item elle est bijective si et seulement si $X$ est \emph{sobre}.
\end{itemize}

Enfin $X \mapsto \widetilde{X}$ est l'adjoint à gauche de l'inclusion des
espaces sobres dans les espaces :
\[
\operatorname{Hom}_{\mathrm{Sp}}(X,\, i Z) \;\simeq\;
\operatorname{Hom}_{\mathrm{Sp\ sob}}(\widetilde{X},\, Z),
\]
c'est-à-dire que $X \to \widetilde{X}$ est l'application universelle de $X$
dans un espace sobre.\footnote{La page porte les deux indices dans l'ordre
inverse, ce que l'adjonction qu'elle énonce juste avant échange ; la
transcription le signale et n'y touche pas. C'est un lapsus de plume : le
membre où figure $\widetilde{X}$ est celui où le but $Z$ est sobre.}

\pagerange{36}{36}
\subsection*{Pourquoi les topos}

« En d'autres termes, $X$ et $\widetilde{X}$ ont la même théorie des
faisceaux » : pour tous les usages pratiques, ou presque, on peut remplacer
l'un par l'autre. De même l'application
\[
\operatorname{Hom}_{\mathrm{Sp}}(X, Y) \longrightarrow
\operatorname{Hom}_{\mathrm{top}}\bigl(\operatorname{Top}(X),
\operatorname{Top}(Y)\bigr)
\]
est bijective dès que $Y$ est sobre, la catégorie de droite étant rigide et
associée à l'ordre de spécialisation.

Puis vient la seule page du dossier où Grothendieck dise ce qu'il cherche, et
elle mérite d'être suivie de près. \emph{Pourquoi la notion de topos est-elle
une généralisation naturelle et utile de celle d'espace ?}

\begin{enumerate}
\item[a)] parce que c'est le cadre naturel de la plupart des structures
algébriques, et de bien des structures non algébriques ;
\item[b)] parce que, à ce titre, elle permet de construire des \emph{topos
classifiants} pour ces structures — peut-être pour toutes : « algèbre $=$
topologie ! » ;\footnote{Le point d'exclamation est de lui. La théorie des
topos classifiants des théories géométriques est aujourd'hui un chapitre
constitué, et cette ligne en est une formulation programmatique : à toute
espèce de structure raisonnable correspond un topos dont les points sont les
modèles.}
\item[c)] parce que la classe des topos est stable par des opérations qui, dans
les espaces, n'auraient pas de sens : le passage au quotient par des relations
d'équivalence sauvages, dont le quotient topologique serait pathologique, et le
recollement de morceaux à la manière d'Artin ;
\item[d)] parce qu'ils apparaissent d'eux-mêmes, comme topos modulaires, dans
diverses questions de géométrie ;
\item[e)] parce qu'ils fournissent le cadre le plus commode partout où l'on
passe systématiquement au quotient ;
\item[f)] parce qu'ils donnent son aboutissement à la notion intuitive de
localisation ;
\item[g)] parce qu'ils sont l'outil technique des théories cohomologiques en
géométrie algébrique.
\end{enumerate}

L'ordre de cette liste est à noter : la géométrie algébrique, qui est le motif
historique, y vient en dernier.

\pagerange{38}{40}
\subsection*{Spécialisations, générateurs, et la question laissée ouverte}

La page 38 revient sur le lien entre $X$ et $\operatorname{Top}(X)$ par les
spécialisations. Si $x \in \overline{\{y\}}$, tout voisinage de $y$ contient
$x$, d'où un homomorphisme fonctoriel de foncteurs fibres
$F_{x} \to F_{y}$, et donc un foncteur de $X$ ordonné par la spécialisation
vers $\mathrm{Fib}(\operatorname{Top}(X))$.

\textbf{Proposition.} \emph{Ce foncteur est pleinement fidèle} :
$\operatorname{Hom}(\xi_{x}, \xi_{y})$ est vide si $x \not\leq y$, et réduit à
l'homomorphisme de spécialisation sinon.

\textbf{Corollaire.} \emph{C'est une équivalence si et seulement si
$X \to \widetilde{X}$ est surjective} — en particulier si $X$ est sobre.

De même, l'ensemble ordonné des applications continues $X \to Y$, pour l'ordre
$f \leq g$ lorsque $f(x) \in \overline{\{g(x)\}}$ pour tout $x$, se plonge
pleinement fidèlement dans $\underline{\operatorname{Hom}}\bigl(
\operatorname{Top}(X), \operatorname{Top}(Y)\bigr)$.\footnote{La page écrit
$\underline{\operatorname{Hom}}(\operatorname{Top}(Y), \operatorname{Top}(X))$,
les deux arguments intervertis. C'est un lapsus, et le manuscrit se corrige
lui-même deux pages plus tôt : la page 36 écrit la même flèche dans le bon
sens, et la page 42 rappellera que $\operatorname{Top}(X)$ dépend de $X$ de
façon covariante.}

Le cours aborde alors la \emph{définition} d'un topos, sous deux formes :
par des propriétés et des générateurs, ou bien par un plongement
\[
E \overset{f_{*}}{\hookrightarrow} \widehat{C},
\qquad C \text{ petite},
\]
pleinement fidèle et admettant un adjoint à gauche $f^{*}$ \emph{exact à
gauche}. Grothendieck note en marge que plusieurs jeux d'axiomes équivalents
sont possibles et se demande, dans une phrase à demi lisible, ce qui gouverne
le choix.

La page 40 précise le second point : on peut prendre pour $C$ n'importe quelle
sous-catégorie génératrice de $E$ et poser $f_{*}(X) = \bigl(A \mapsto
\operatorname{Hom}(A,X)\bigr)$ ; dire que ce foncteur est \emph{conservatif}
signifie que $C$ est génératrice, et cela équivaut à la pleine fidélité de
$f_{*}$, laquelle entraîne l'existence de l'adjoint à
gauche.\footnote{L'équivalence entre « conservatif » et « pleinement fidèle »
n'est pas formelle : elle vaut dans un topos, où les épimorphismes sont
effectifs et les colimites universelles, et c'est ce qui fait qu'une famille
génératrice y est automatiquement dense. La page hésite précisément là :
« does imply » est barré et remplacé par « is equivalent with ». La correction
va dans le bon sens.} Suivent les usages des générateurs — existence
d'adjoints, représentabilité des foncteurs contravariants transformant les
colimites en limites — puis les exemples : le topos vide
$\operatorname{Top}(\emptyset)$ et le topos ponctuel $\operatorname{Top}(1)$ ;
$\operatorname{Top}(X,G)$, avec $\operatorname{Top}(*,G) = B_{G}$ et
$H^{i}(B_{G}, -) = H^{i}(G,-)$, le foncteur fibre de $B_{G}$ étant le foncteur
d'oubli ; et les $\widehat{C}$, avec le plongement pleinement fidèle
$C^{\circ} \to \mathrm{Fib}(\widehat{C})$, qui est même une équivalence
$\mathrm{Ind}(C^{\circ}) \simeq \mathrm{Fib}(\widehat{C})$.

La page s'achève sur une question, et c'est elle qui court d'un bout à l'autre
du dossier :

\begin{quote}
Peut-on avoir une équivalence $\widehat{C} \simeq \operatorname{Top}(X)$ ?
\end{quote}

Il en esquisse aussitôt la réponse : il suffirait que $C$ soit une catégorie
ordonnée et que $X$ ait des points admettant un plus petit voisinage — des
points essentiels — qui soient denses ; on prendrait alors $X = \mathrm{Ind}(C)$.
La page s'arrête là, plusieurs mots illisibles.\footnote{La réponse complète,
dans les deux sens et avec ses conditions exactes, est à la page 48,
qui reprend la question dans les mêmes termes et dans la même langue. Elle
n'est pas ailleurs dans le dossier : c'est la seule continuité de cette
portée dans la moitié anglaise du cours.}

\pagerange{42}{42}
\section*{Morphismes de topos (pages 42, 44, 46)}

L'autre manuscrit, celui des pages paires, reprend son cours. Il est ici en
français ; les pages 48, 52, 53 et 54 seront en anglais, et l'on en donne comme
plus haut une restatement française et non une traduction.

\pagerange{42}{42}
\subsection*{Les topos forment une 2-catégorie}

Un morphisme de topos $f : X \to Y$ est déterminé par son image directe, le
foncteur
\[
\underline{\mathrm{Mor\,Top}}(X,Y) \hookrightarrow
\underline{\operatorname{Hom}}_{\mathrm{cat}}(X,Y), \qquad f \mapsto f_{*},
\]
étant pleinement fidèle ; la composition en fait une 2-catégorie. La notion
d'\emph{équivalence} de topos — $f$ tel qu'il existe $g$ avec $gf \simeq
\mathrm{id}$ et $fg \simeq \mathrm{id}$ — y remplace exactement la notion
d'homéomorphisme pour les espaces, et c'est l'équivalence usuelle de
catégories. La page prend soin de le noter en marge : il n'y a plus
d'isomorphisme, seulement de l'équivalence.

\pagerange{42}{44}
\subsection*{Les exemples, un à un}

\begin{itemize}
\item \emph{Les espaces.} $\operatorname{Top}(X)$ dépend de $X$ de façon
covariante, et c'est un foncteur 2-fidèle des espaces vers les topos ; pour
$\operatorname{Top}(X,G)$ la pleine fidélité tombe.

\item \emph{Morphismes vers un topos spatial.} Se donner $E \to
\operatorname{Top}(X)$, c'est se donner une application
\[
\mathrm{Ouv}(X) \longrightarrow \{\text{sous-objets de } e_{E}\}
\]
commutant aux bornes supérieures quelconques et aux intersections finies —
c'est-à-dire un morphisme de cadres, c'est-à-dire un morphisme de
locales.\footnote{C'est l'énoncé qui fait des topos localiques la partie
« spatiale » de la théorie, et c'est la même donnée que celle du feuillet 28,
où $\gamma$ était exactement une telle application.}

\item \emph{Le doublet.} Soit $D = \{x,y\}$ l'espace à deux points dont l'un est
fermé et l'autre générique. Un morphisme $E \to \operatorname{Top}(D)$ est
exactement la donnée d'un ouvert de $E$ : $\operatorname{Top}(D)$ est le
\emph{topos classifiant} de l'espèce de structure « sous-objet de l'objet
final ». Et tout espace topologique se plonge dans un produit $D^{I}$, avec
$I = \mathrm{Ouv}(X)$.\footnote{C'est le premier topos classifiant explicite du
dossier, et le plus simple qui soit : l'espace de Sierpiński classifie les
valeurs de vérité. La page 36 annonçait le programme — « algèbre $=$
topologie ! » — ; en voici le cas de base.}

\item \emph{La réflexion localique.} Si l'ensemble ordonné des sous-objets de
$e_{E}$ est isomorphe à $\mathrm{Ouv}(X)$ pour un espace $X$, alors le
morphisme $E \to \operatorname{Top}(X)$ ainsi obtenu est \emph{universel} parmi
les morphismes de $E$ vers un espace.

\item \emph{Les points.} Les morphismes d'un point dans $E$ forment la
catégorie $\underline{\mathrm{Pt}}(E)^{\circ}$.

\item \emph{Les topos de préfaisceaux.} Un foncteur $u : C \to C'$ donne la
restriction $u^{*} : \widehat{C}' \to \widehat{C}$, laquelle possède un adjoint
à gauche $u_{!}$ et un adjoint à droite $u_{*}$ — les deux extensions de
Kan :
\[
u_{!} \dashv u^{*} \dashv u_{*}.
\]
\footnote{La page note $u^{!}$ l'adjoint à droite et l'attribution des
propriétés de commutation y est brouillée ; nous écrivons $u_{*}$, notation
constante ailleurs dans le dossier, et rétablissons la chaîne. C'est la
possession de \emph{deux} adjoints qui fait qu'un morphisme entre topos de
préfaisceaux issu d'un foncteur de sites est essentiel, au sens de la page 15.}

\item \emph{Sortir d'un topos de préfaisceaux.} Pour tout topos $E$,
\[
\operatorname{Hom}_{\mathrm{top}}(\widehat{C}, E) \;\simeq\;
\operatorname{Hom}(C, \underline{\mathrm{Pt}}(E)),
\]
et en particulier
$\operatorname{Hom}_{\mathrm{top}}(\widehat{C}, \widehat{C}') \simeq
\operatorname{Hom}(C, \mathrm{Pro}\,C')$, le sous-ensemble
$\operatorname{Hom}(C,C')$ y étant plongé pleinement fidèlement. Et si $C$ est
stable par limites finies, on retrouve dans l'autre sens la formule de la page 6 :
$\operatorname{Hom}_{\mathrm{top}}(E, \widehat{C}) \simeq
\operatorname{Hom}_{\mathrm{ex.g.}}(C, E)^{\circ}$.\footnote{Le second membre
demande la même convention de 2-flèches que celle discutée à la page 6 :
c'est elle que note l'exposant $(\ )^{\circ}$. De même
$\underline{\mathrm{Pt}}(\widehat{C}') \simeq \mathrm{Pro}(C')$ est la forme
opposée de l'énoncé $\mathrm{Ind}(C^{\circ}) \simeq \mathrm{Fib}(\widehat{C})$
de la page 40.}
\end{itemize}

\pagerange{46}{46}
\subsection*{Les cas dégénérés, et l'induction}

La page 46 achève l'énumération. Vers le \emph{topos ponctuel}, l'image inverse
donne les objets constants et l'image directe le foncteur des sections
globales. Quant au \emph{topos vide}, un morphisme $E \to
\operatorname{Top}(\emptyset)$ n'existe que si $E \simeq
\operatorname{Top}(\emptyset)$ lui-même : l'image inverse devrait envoyer
l'unique objet à la fois sur l'objet initial et sur l'objet final de $E$, ce
qui les identifie et rend $E$ ponctuel.

Vient enfin le morphisme d'\emph{induction} $j_{X} : E_{/X} \to E$, qui portera
tout ce qui suit :
\[
j_{X}^{*}(Y) = Y \times X, \qquad
j_{X!}(Y \to X) = Y,
\]
le troisième, $j_{X*}$, étant « plus compliqué » de l'aveu de la page, et
n'y étant pas explicité. On a $j_{X!} \dashv j_{X}^{*} \dashv j_{X*}$ — d'où le fait, que la page
souligne, que $j_{X}^{*}$ commute aux limites \emph{et} aux colimites
quelconques. Comme foncteur en $X$,
\[
\operatorname{Hom}_{\mathrm{top}_{E}}(E_{/X}, E_{/Y}) \;\simeq\;
\operatorname{Hom}_{E}(X,Y) \quad \text{(catégorie discrète)}.
\]
Un dernier titre, « topos somme de topos », reste sans texte.

\pagerange{48}{48}
\section*{Quand un topos de préfaisceaux est-il spatial ? (page 48)}

C'est la réponse à la question laissée ouverte au bas de la page 40, et il
faut la lire d'affilée avec elle : deux pages éloignées, une seule question.

\textbf{Sens direct.} \emph{$\widehat{C}$ est équivalent à un
$\operatorname{Top}(X)$ si et seulement si $C$ est une catégorie ordonnée},
c'est-à-dire si $x \to e_{\widehat{C}}$ est un monomorphisme pour tout objet
$x$. On prend alors pour $X$ l'ensemble $\mathrm{Ob}\,C$, muni de la topologie
engendrée par les $O_{x} = \{ y \mid y \leq x \}$, lesquels en forment une
\emph{base} ; les ouverts sont les cribles.\footnote{La page 48 écrit l'ordre à
l'envers de celui fixé à la page 23 : ici $O_{x}$ est l'ensemble des
$y \leq x$, là celui des $y \geq x$. Les deux énoncés — $\operatorname{Top}(X)
\simeq \operatorname{Hom}(\underline{X}, \mathrm{Ens})$ à la page 24, et celui-ci —
sont le même une fois la traduction faite. Nous gardons partout la convention
de la page 23 et signalons ici l'échange plutôt que de renuméroter la page.}

La démonstration tient en une observation : \emph{les recouvrements sont
triviaux}. Si les $O_{x_{i}}$ recouvrent $O_{x}$, l'un d'eux contient $x$,
donc $x \leq x_{i}$ ; mais $x_{i} \leq x$ puisque $O_{x_{i}} \subset O_{x}$,
d'où $x \approx x_{i}$ et $O_{x_{i}} = O_{x}$. Un faisceau sur cette base n'est
donc astreint à rien de plus qu'un préfaisceau sur $C$, et l'on a bien
$\operatorname{Top}(X) \simeq \widehat{C}$.

\textbf{Sens réciproque.} \emph{Si $X$ est un espace sobre,
$\operatorname{Top}(X)$ est équivalent à un $\widehat{C}$ si et seulement si
les points essentiels de $X$ — ceux qui ont un plus petit voisinage ouvert, ou,
ce qui revient au même, ceux dont le foncteur fibre commute aux produits
quelconques — sont très denses}, et l'on peut prendre pour $C$ la catégorie de
ces points.

Cet énoncé referme exactement la liste de conditions équivalentes de la
page 26, et lui donne sa forme constructive : la page 26 disait qu'un topos est
à la fois spatial et de préfaisceaux si et seulement s'il est engendré par les
sous-objets de son objet final et a suffisamment de points essentiels ; la
page 48 dit comment fabriquer $X$ à partir de $C$, et $C$ à partir de
$X$.\footnote{La caractérisation d'un point essentiel par la commutation du
foncteur fibre aux produits quelconques n'apparaît que sur cette page. C'est
elle qui explique la terminologie des pages 23 à 26 : un point est essentiel exactement
quand le morphisme de topos qu'il définit est essentiel, c'est-à-dire quand son
image inverse a elle-même un adjoint à gauche.}

\pagerange{49}{49}
\section*{Images directes et inverses, localisation, Galois (pages 49 à 51)}

\pagerange{49}{49}
\subsection*{Les trois foncteurs et leurs prolongements}

Soit $f : C \to C'$ un foncteur, $g$ son adjoint à droite et $h$ celui de $g$,
lorsqu'ils existent. Les trois foncteurs entre topos de préfaisceaux se
répartissent alors ainsi : $f_{!}$ prolonge $f$, $f^{*}$ prolonge $g$, $f_{*}$
prolonge $h$. Si $f$ commute aux limites projectives finies, $f_{!}$ y commute
encore, et l'on obtient un morphisme de topos $\varphi : \widehat{C}' \to
\widehat{C}$ avec
\[
\varphi^{*} = f_{!}, \qquad \varphi_{*} = f^{*}.
\]
\footnote{Que $f_{!}$, extension de Kan à gauche le long de Yoneda, soit exacte
à gauche lorsque $f$ l'est et que $C$ a des limites finies, est le fait qu'on
attache aujourd'hui au nom de Diaconescu. La page l'énonce sans le démontrer et
sans commentaire.}

Les deux calculs de la page sont les identités d'adjonction lues sur les
représentables : $f_{!}(x) = f(x)$ d'une part, et de l'autre
\[
f_{*}(x) = \bigl(x' \longmapsto \operatorname{Hom}_{\widehat{C}}
(x' \circ f^{\circ},\, x)\bigr)
= \bigl(x' \longmapsto \operatorname{Hom}_{C}(g(x'),\, x)\bigr)
\]
lorsque $g$ existe. La page porte enfin, à part, « $(\mathbb{Z}, G)$ » et
« $B_{G} \to B_{e}$ » : c'est l'annonce des deux pages suivantes.

\pagerange{50}{50}
\subsection*{Un sous-groupe, c'est une localisation}

Voici le morceau le plus frappant du dossier. Soient $G$ un groupe d'un topos
$E$, $B_{G}$ son topos classifiant, $H \subset G$ un sous-groupe et
$X = G/H$, muni de son point marqué $\varepsilon \in \Gamma(X)$. On a alors un
carré

\begin{tikzcd}[row sep=large]
(B_{G})_{/X} \arrow[r, "\varphi"] \arrow[d, "j_{X}"'] & B_{H} \arrow[d, "B_{i}"] \\
B_{G} & B_{G}
\end{tikzcd}

dans lequel $\varphi$ est une \emph{équivalence}, donnée dans les deux sens par
\[
\varphi(E') = E' \times_{X} (e, \varepsilon), \qquad
\psi(F) = G \wedge_{H} F,
\]
et l'on vérifie $\varphi(j_{X}^{*}(E')) = (E' \times X) \times_{X}
(e, \varepsilon) \simeq E'$.\footnote{Le carré du manuscrit porte $B_{G}$ deux
fois en bas, sans flèche entre les deux sommets : c'est le même topos, et
c'est cette identification que le raisonnement utilise. La marge glose
$\varphi$ d'un mot juste : « oubli partiel des opérations de $G$ — ou plutôt de
$H$ ».}

Le sens de tout cela est que \emph{se restreindre à un sous-groupe et se
localiser au-dessus de $G/H$ sont la même opération}. Via $\varphi$, le
morphisme d'induction $j_{X}$ s'interprète comme $B_{i}$, et les trois foncteurs
deviennent les trois opérations classiques de la théorie des représentations :
\[
j_{X}^{*} = \text{restriction}, \qquad
j_{X!} : F \mapsto G \wedge_{H} F \ (\text{induction}), \qquad
j_{X*} : F \mapsto \operatorname{Hom}_{H}(G, F) \ (\text{coinduction}).
\]

\pagerange{50}{51}
\subsection*{Et c'est la théorie de Galois}

La page tire aussitôt le cas des groupes discrets, et c'est là que l'énoncé
formel devient une théorie connue. Si $G = \pi_{1}(Z)$ est le groupe
fondamental d'un espace $Z$ connexe, localement connexe et localement
1-connexe, alors $B_{G}$ est la catégorie des revêtements de $Z$ ; la donnée de
$X = G/H$ est celle d'un revêtement pointé au-dessus du point base ; et
l'équivalence $\varphi$ ci-dessus devient l'équivalence
\[
\mathrm{Rev}(X) \;\simeq\; B_{H}, \qquad H \simeq \pi_{1}(X,x),
\]
c'est-à-dire \emph{la théorie de Galois du revêtement $X$}. Le dictionnaire
entre sous-groupes du groupe fondamental et revêtements tombe comme cas
particulier d'une remarque sur les topos induits.\footnote{La marge porte
« théorie de Galois pour $Z$ locaux », lecture incertaine. C'est bien la
correspondance galoisienne des revêtements, dont la version schématique — le
groupe fondamental comme groupe d'automorphismes d'un foncteur fibre — est le
sujet de SGA 1 V. Le mérite de cette demi-page est de la présenter non comme un
théorème mais comme une lecture : un revêtement \emph{est} une localisation du
topos classifiant.}

La page 51 refait le même calcul avec un espace : si un groupe discret $G$ opère
sur un espace $Z$, le foncteur « restriction des opérations à $H$ »
\[
\operatorname{Top}(Z, G) \longrightarrow \operatorname{Top}(Z, H)
\]
est encore un foncteur de localisation. On pose $Z' = Z \times G/H$ avec
l'action diagonale, et le foncteur
$\operatorname{Top}(Z,G)_{/Z'} \to \operatorname{Top}(Z,H)$, $E' \mapsto E'
\times_{Z'} (X,e)$, est une équivalence — $Z'$ étant, comme le note la page, un
« revêtement topologique » de $Z$.\footnote{La page note $\mathbb{Z}$ l'espace
sur lequel le groupe opère, là où $Z$ ou $X$ étaient attendus, et la phrase
finale est laissée bancale (« on le foncteur \ldots\ est une équivalence »).
Nous écrivons $Z$ partout.}

\pagerange{52}{52}
\section*{Sites (pages 52, 54, 53)}

Les trois dernières pages sont un condensé de la théorie des sites, en anglais.
\emph{Elles sont reliées dans le désordre} : la page 54 précède la page 53, et
c'est dans cet ordre qu'on les lit ici.\footnote{Deux indices concordants, et
ils suffisent. La numérotation des théorèmes : Th. 1 et Th. 2 sont sur la
page 54, Th. 3, Th. 5 et Th. 4 sur la page 53. Et la phrase : la page 54
s'interrompt sur « \emph{it is the finest topology for which} », que la
première ligne de la page 53 achève exactement — « \emph{the objects of
$\widetilde{C}$ [are] sheaves} ». L'ordre 52, 54, 53 est donc le nôtre, mais il
n'est pas une conjecture : il est ce que les deux feuillets disent d'eux-mêmes.
Il s'ensuit que le dossier ne s'arrête pas au milieu d'une phrase, contrairement
à ce que l'ordre matériel donne à croire : il s'arrête après le lemme de
comparaison, au bas de la page 53. Nous conservons par ailleurs la numérotation
du manuscrit, y compris l'ordre Th. 3, Th. 5, Th. 4 de la page 53 : c'est sous
ces numéros que la page renvoie à ses propres énoncés, et les renuméroter
rendrait le renvoi faux.}

\pagerange{52}{52}
\subsection*{Ce qu'est une topologie}

Le plan est celui d'un cours : (1) la notion de localisation, « yoga » ; (2) la
donnée d'une localisation sur une catégorie, c'est-à-dire une
\emph{topologie}, sous deux formes équivalentes.

\emph{En termes de familles couvrantes}, lorsque les produits fibrés existent :
stabilité par changement de base (Cov 1), par composition (Cov 2), présence des
identités (Cov 3), et, entre crochets, la saturation (Cov 4). La page note en
accolade qu'il suffit de supposer les $X_{\alpha} \to X$ carrables.

\emph{En termes de cribles couvrants} : changement de base (T1), caractère
local (T2), et une troisième condition biffée. Les topologies sont alors
ordonnées par l'inclusion, se coupent, ont des bornes supérieures, et sont
engendrées par des familles de cribles ou de morphismes.

Le lien entre les deux formes est l'identité
\[
\operatorname{Hom}_{\widehat{C}}(R, F) \;=\; \varprojlim_{C/R} F(X)
\;\simeq\; \operatorname{Ker}\Bigl( \prod_{\alpha} F(X_{\alpha})
\rightrightarrows \prod_{\alpha,\beta} F(X_{\alpha} \times_{X} X_{\beta})
\Bigr),
\]
qui dit qu'un crible engendré par une famille impose exactement la condition de
recollement de cette famille.\footnote{La page omet le premier produit et écrit
$\operatorname{Ker}(F(X_{\alpha}) \rightrightarrows \ldots)$ ; c'est une
abréviation, non une erreur de fond.} Sous Cov 1 à Cov 3, les cribles engendrés
par les familles couvrantes sont \emph{cofinaux} parmi les cribles couvrants —
c'est ce qui rend les deux présentations interchangeables.

Vient (3) la notion de faisceau, et celle de préfaisceau séparé, données d'un
seul geste : $F$ est un faisceau (resp. séparé) lorsque
$\operatorname{Hom}(X,F) \to \operatorname{Hom}(R,F)$ est bijective (resp.
injective) pour tout crible couvrant $R \subset X$. Enfin, étant donnée une
famille $\Phi$ de préfaisceaux, il existe une \emph{plus fine} topologie
$T_{\Phi}$ faisant de tout $F \in \Phi$ un faisceau, décrite explicitement par
changement de base.

\pagerange{54}{54}
\subsection*{Le faisceau associé}

La page 54 construit le foncteur $L$ :
\[
LF(X) \;=\; \varinjlim_{R \in T(X)} \operatorname{Hom}(R, F),
\qquad \ell : \mathrm{id} \longrightarrow L,
\]
avec les quatre propriétés qui font toute la construction : a) $L$ est exact à
gauche ; b) $LF$ est toujours séparé ; c) $F$ est séparé si et seulement si
$F \to LF$ est un monomorphisme, et alors $LF$ est un faisceau ; d) $F \to LF$
est un isomorphisme si et seulement si $F$ est un faisceau.

\textbf{Théorème 1.} \emph{L'inclusion $\widetilde{C} \subset \widehat{C}$
admet un adjoint à gauche $a$, et $a$ est exact à gauche} — car $a = L \circ L$,
b) et c) donnant la conclusion en deux applications.\footnote{L'exactitude à
gauche de $L$ repose sur le fait que $T(X)$ est cofiltrant pour l'inclusion des
cribles ; la page ne le mentionne pas. C'est la même hypothèse qui rend la
colimite définissant $LF$ exacte à gauche.}

\textbf{Corollaires.} Les limites projectives existent dans $\widetilde{C}$ et
l'inclusion y commute ; les colimites existent, données par
$\varinjlim^{\widetilde{C}} F_{i} \simeq a\bigl(\varinjlim^{\widehat{C}}
F_{i}\bigr)$ ; et tout faisceau est colimite de ses représentables,
$F \simeq \varinjlim_{C/F} a(X)$.\footnote{La page glose, et la glose vaut
d'être gardée : de là suivent « toutes les propriétés d'exactitude de
$(\mathrm{Ens})$ qui n'engagent que des limites finies et des colimites
quelconques ». C'est exactement la clause qui fait d'un topos une catégorie où
l'on calcule comme dans les ensembles. Le corollaire sur les colimites est
d'ailleurs celui qu'utilisait la page 20.}

\textbf{Théorème 2.} \emph{Soit $\varepsilon_{C} : C \to \widetilde{C}$,
$X \mapsto a(X)$, qui est exact à gauche. Alors} a) \emph{une famille
$X_{\alpha} \to X$ est couvrante si et seulement si
$\coprod \varepsilon_{C}(X_{\alpha}) \to \varepsilon_{C}(X)$ est un
épimorphisme ;} b) \emph{un crible $R \subset X$ est couvrant si et seulement
si $\varepsilon_{C}(R) \xrightarrow{\ \sim\ } \varepsilon_{C}(X)$.}\footnote{L'exactitude
à gauche de $\varepsilon_{C}$ suppose que $C$ ait des limites finies ; sans
cela, $\varepsilon_{C}$ est seulement plat. La page ne le dit pas ici, mais le
suppose partout ailleurs dans le dossier.}

On récupère ainsi la topologie $T$ à partir de la seule donnée de
$\widetilde{C} \subset \widehat{C}$ : \emph{c'est la topologie la plus fine
pour laquelle les objets de $\widetilde{C}$ sont des faisceaux.}\footnote{C'est
la phrase que la page 54 laisse en suspens et que la page 53 achève ; le mot
« finest » est donné pour une lecture douteuse par la transcription, et il est
bien le bon : les topologies faisant de tout objet de $\widetilde{C}$ un
faisceau sont, par le Th. 3 ci-dessous qui renverse l'ordre, exactement celles
qui sont moins fines que $T$ — d'où $T$ comme plus fine d'entre elles. La
page 52 emploie d'ailleurs le même mot dans la même construction.}

\pagerange{53}{53}
\subsection*{Les théorèmes de Giraud}

\textbf{Théorème 3} (Giraud). \emph{$T \mapsto \widetilde{C}_{T}$ est une
correspondance bijective entre les topologies sur $C$ et les sous-catégories
strictement pleines $E \subset \widehat{C}$ dont l'inclusion admet un adjoint à
gauche exact à gauche. Elle renverse l'ordre :
$T \leq T' \iff \widetilde{C}_{T'} \subset \widetilde{C}_{T}$.}\footnote{C'est
la correspondance entre topologies de Grothendieck sur $C$ et localisations
exactes à gauche de $\widehat{C}$ — c'est-à-dire sous-topos de $\widehat{C}$ —
dont la formulation interne, par les topologies de Lawvere--Tierney, est
postérieure. Elle est le pendant, du côté des sites, de l'énoncé sur les
sous-topos dont la page 3 se servait pour prouver que $\Pi(E,F)$ est un topos.}

\textbf{Théorème 5.} \emph{Une catégorie $E$ est équivalente à un
$\widetilde{C}$, pour une petite $C$, si et seulement si elle vérifie les
axiomes} a) b) c) d)\footnote{Les quatre axiomes ne sont pas écrits sur la
page — le cours devait les avoir donnés ailleurs — et nous ne les suppléons
pas : ce sont les axiomes de Giraud, et les nommer suffit à situer l'énoncé
sans lui prêter une formulation qui n'est pas la sienne.}. \emph{On peut alors
prendre pour $C$ n'importe quelle sous-catégorie pleine génératrice de $E$,
munie de la topologie induite par la topologie canonique de $E$ ; et la donnée
d'une équivalence $\widetilde{C}_{T} \simeq E$ équivaut à celle d'un foncteur
pleinement fidèle $C \to E$ d'image génératrice.}

\textbf{Théorème 4} (lemme de comparaison). \emph{Soit $C \hookrightarrow C'$
un foncteur pleinement fidèle entre sites, la topologie de $C$ étant induite. Si
tout objet de $C'$ peut être couvert par des objets de $C$, alors
$\widetilde{C}' \xrightarrow{\ \sim\ } \widetilde{C}$.} La page ajoute qu'il y a
une réciproque lorsque la topologie de $C$ est moins fine que la topologie
canonique.

La marge porte les derniers mots du dossier, et ils sont d'un praticien :
« analogue bien connu en topologie ordinaire ; possibilité de définir des
$\widetilde{C}$ équivalents de bien des façons différentes ». C'est là que la
chemise s'arrête.\footnote{Dans l'ordre matériel des feuillets, la dernière
page est la 54 et le dossier paraît s'interrompre au milieu d'une phrase. Dans
l'ordre que les deux pages indiquent elles-mêmes, il s'achève ici, sur une
remarque de méthode et non sur une rupture. Nous n'ajoutons rien à l'une ni à
l'autre.}

\end{document}
